%%%%%%%%%%%%%%%%%%%%%%%%%%%%%%%%%%%%%%%%%%%%%%%%%%%%%%%%%%%%%%%%%%%%%%
% Lớp 11 | Dãy số. Cấp số cộng. Cấp số nhân | Cấp số nhân
% Dạng 1D2?3-4: Tìm số hạng cụ thể trong cấp số nhân
% Số câu: 114  (sinh tự động bởi tach_cau.py; từ nay thầy sửa trực tiếp file này)
%%%%%%%%%%%%%%%%%%%%%%%%%%%%%%%%%%%%%%%%%%%%%%%%%%%%%%%%%%%%%%%%%%%%%%

%%%%%%%%%%%%%%%% TRẮC NGHIỆM 4 PHƯƠNG ÁN %%%%%%%%%%%%%%%%

% nguồn: Ok_Toan11_TracNghiem_Lan2
\begin{ex}%[1D2N3-4]%[ma:47e4261b]
	Cho cấp số nhân có $u_1=-3$, $q=\dfrac{2}{3}$.  Tính $u_5$.
	\choice 
		 {$u_5=\dfrac{-27}{16}$}
		{\True  $u_5=\dfrac{-16}{27}$}
		{$u_5=\dfrac{16}{27}$} 
		{$u_5=\dfrac{27}{16}$} 
	\loigiai{Ta có:$u_5 = u_1 q^4 =(-3) \left(\dfrac{2}{3} \right)^4 = \dfrac{-16}{27}$.} 
\end{ex}

% nguồn: Ok_Toan11_TracNghiem_Lan2
\begin{ex}%[1D2N3-4]%[ma:e0ea1ce9]
	Cho cấp số nhân $(u_n)$,  biết $u_1=12$, $\dfrac{u_3}{u_8}=243$. Tìm $u_9$.
	\choice 
		{$u_9=\dfrac{2}{2\,187}$} 
		 {$u_9=\dfrac{4}{6\,563}$}
		{$u_9=78\,732$} 
		 {\True $u_9=\dfrac{4}{2\,187}$}
	\loigiai{Gọi $q$ là công bội của  cấp số nhân $(u_n)$. Ta có
		\[u_3=u_1q^2,~u_8=u_1q^7 \Rightarrow \dfrac{u_3}{u_8}=\dfrac{1}{q^5}=243 \Rightarrow q=\dfrac{1}{3}\] 
	Do đó  $u_9=u_1q^8=12\left(\dfrac{1}{3}\right)^8=\dfrac{4}{2\,187}$.}
\end{ex}

% nguồn: Ok_Toan11_TracNghiem_Lan2
\begin{ex}%[1D2N3-4]%[ma:7514a9c5]
	Cho dãy số $\left(u_n \right)$ xác định bởi  $\heva{&u_1=1\\&u_{n+1}=2u_n+5}$. Tính số hạng thứ $2\,018$ của dãy số trên
	 \choice
		{\True $u_{2\,018}=6\cdot 2^{2\,017}-5$} 
		{$u_{2\,018}=6\cdot  2^{2\,018}-5$}
		{$u_{2\,018}=6\cdot 2^{2\,017}+1$} 
		{$u_{2\,018}=6\cdot  2^{2\,018}+5$}
	\loigiai{Đặt $v_n=u_n+5$. Ta có 
		 $v_{n+1}=u_{n+1}+5=2u_n+5+5=2(u_n+5)=2v_n$.\\
		Do đó $v_n$ là cấp số nhân với $v_1=6, q=2$,  $v_n=6\cdot 2^{n-1}$\\
		 $v_{2\,018}=6\cdot 2^{2\,017}
	\Rightarrow  u_{2\,018}=6\cdot 2^{2\,017}-5$.}
\end{ex}

% nguồn: Ok_Toan11_TracNghiem_Lan2
\begin{ex}%[1D2N3-4]%[ma:7d2e86ee]
	Cho cấp số nhân có các số  hạng lần lượt là $1$; $4$; $16$; $64$; $\cdots$. Gọi $S_n$ là tổng của $n$ số hạng đầu tiên của cấp số nhân đó. Mệnh đề nào sau đây đúng?
	\choice 
		{$S_n=4^{n-1}$} 
		{$S_n=\dfrac{n\left(1+4^{n-1}\right)}{2}$} 
		{\True  $S_n=\dfrac{4^n-1}{3}$}
		{$S_n=\dfrac{4\left(4^n-1\right)}{3}$} 
	\loigiai{ 
		Cấp số nhân  đã cho có $\heva{&u_1=1\\&q=4.}$\\
		Suy ra  $S_n=u_1\cdot \dfrac{1-q^n}{1-q}=1\cdot \dfrac{1-4^n}{1-4}=\dfrac{4^n-1}{3}$.
	} 
\end{ex}

% nguồn: Ok_Toan11_TracNghiem_Lan2
\begin{ex}%[1D2N3-4]%[ma:b26a1ec7]
	Cho cấp số nhân $(u_n)$ có $u_1=5$, công bội $q=-2$. Số  hạng thứ sáu của $(u_n)$ là
	\choice 
		 {$u_6=320$}
		 {\True $u_6=-160$}
		{$u_6=-320$} 
		{$u_6=160$} 
	 \loigiai{
	Ta có $u_6=u_1q^5=5\cdot (-2)^5=-160$. 
	} 
\end{ex}

% nguồn: Ok_Toan11_TracNghiem_Lan2
\begin{ex}%[1D2N3-4]%[ma:d188d0b6]
	 Tìm số hạng đầu $u_1$ của cấp số nhân $\left(u_n \right)$ biết rằng $u_1+u_2+u_3=168$ và $u_4+u_5+u_6=21$
	\choice 
		{$u_1=24$} 
		 {$u_1=\dfrac{1334}{11}$}
		 {\True $u_1=96$}
		{$u_1=\dfrac{217}{3}$} 
	\loigiai{Ta có 
		\begin{eqnarray*}
			& &\begin{cases} u_1+u_2+u_3=168 \\ u_4+u_5+u_6=21
		\end{cases}\\
		&\Leftrightarrow &\begin{cases} u_1+u_1q+u_1q^2=168 \\ u_1q^3+u_1q^4+u_1q^5=21
	\end{cases}\\
	&\Leftrightarrow &\begin{cases} u_1\left(1+q+q^2\right) =168 \\ u_1q^3\left(1+q+q^2\right) =21
	\end{cases}\\
	&\Leftrightarrow &\begin{cases} u_1=\dfrac{168}{1+q+q^2} \\ q^3=\dfrac{21}{168}
	\end{cases}\\
	&\Leftrightarrow &\begin{cases} u_1=96 \\ q^3=\dfrac{1}{8}
	\end{cases}\\
	&\Leftrightarrow &\begin{cases} u_1=96 \\ q=\dfrac{1}{2}.
	\end{cases}
	\end{eqnarray*}
Vậy $u_1=96$.} 
\end{ex}

% nguồn: Ok_Toan11_TracNghiem_Lan2
\begin{ex}%[1D2N3-4]%[ma:6905d522]
	Cho cấp số nhân $(u_n)$ có $u_1=-2$ và công bội $q=3$. Số  hạng $u_2$ là
	\choice 
		{\True $u_2=-6$} 
		{$u_2=6$} 
		 {$u_2=1$}
		{$u_2=-18$} 
	\loigiai{ 
	Ta có  $u_2=u_1q=-2\cdot 3=-6$.
	 }
\end{ex}

% nguồn: Ok_Toan11_TracNghiem_Lan2
\begin{ex}%[1D2N3-4]%[ma:fb36ba9f]
	Cho cấp số nhân  có $u_1=-3$, $q=\dfrac{2}{3}$. Số $\dfrac{-96}{243}$ là số hạng thứ mấy của cấp số này?
	 \choice
		{Thứ  $5$}
		{Thứ $7$} 
		{\True Thứ $6$} 
		{Không phải là số hạng của cấp  số}
	 \loigiai{
		Giả sử số $\dfrac{-96}{243}$ là hạng thứ $n$ của cấp số  này.\\
		 Ta có
		\[u_1 q^{n-1} = \dfrac{-96}{243} \Leftrightarrow(-3) \left(\dfrac{2}{3} \right)^{n-1} = \dfrac{-96}{243} \Leftrightarrow n =  6.\]
	Vậy số  $\dfrac{-96}{243}$ là số hạng thứ $6$ của cấp số.}
\end{ex}

% nguồn: Ok_Toan11_TracNghiem_Lan2
\begin{ex}%[1D2N3-4]%[ma:e9571936]
	Cho cấp số nhân $(u_n)$ có  $u_1=3$, công bội $q=2$. Tính $u_5$.
	 \choice
		 {$24$}
		{$11$} 
		 {\True $48$}
		{$9$} 
	 \loigiai{
	Ta  có $u_5=u_1q^4=3\cdot 2^4=48$.
	 }
\end{ex}

% nguồn: Ok_Toan11_TracNghiem_Lan2
\begin{ex}%[1D2N3-4]%[ma:e712c862]
	Cho cấp số nhân $(u_n)$ có công bội $q>0$. Biết  $u_1=1$, $u_3=4$. Tìm $u_4$.
	 \choice
		 {$\dfrac{11}{2}$}
		{$2$} 
		 {$16$}
		 {\True $8$}
	 \loigiai{
	Ta có $u_3=u_1q^2\Leftrightarrow  4=1\cdot q^2\Rightarrow q=2$ (vì $q>0$).\\
	 Vậy $u_4=u_1q^3=8$.
	 }
\end{ex}

% nguồn: Ok_Toan11_TracNghiem_Lan2
\begin{ex}%[1D2N3-4]%[ma:76130976]
	Cho cấp số nhân có $u_1=-3$ và công bội  $q=\dfrac{2}{3}$. Tính $u_5$.
	 \choice
		{$u_5=-\dfrac{27}{16}$} 
		 {\True $u_5=-\dfrac{16}{27}$}
		{$u_5=\dfrac{16}{27}$} 
		 {$u_5=\dfrac{27}{16}$}
	\loigiai{ 
	Ta có  $u_5=u_1q^4=(-3)\cdot\left(\dfrac{2}{3}\right)^4=-\dfrac{16}{27}$.
	 }
\end{ex}

% nguồn: Ok_Toan11_TracNghiem_Lan2
\begin{ex}%[1D2N3-4]%[ma:491bd0c7]
	Cho cấp số nhân $(u_n)$ có $u_1=2$,  công bội $q=3$. Giá trị của $u_{2\,019}$ bằng
	\choice 
		{\True $2\cdot  3^{2\,018}$}
		 {$3\cdot 2^{2\,018}$}
		{$2\cdot  3^{2\,019}$}
		{$3\cdot  2^{2\,019}$}
	 \loigiai{
	Ta có $u_{2\,019}=u_1q^{2\,018}=2\cdot 3^{2\,018}$. 
	} 
\end{ex}

% nguồn: Ok_Toan11_TracNghiem_Lan2
\begin{ex}%[1D2N3-4]%[ma:39687d32]
	Cho cấp số nhân $(u_n)$ có $u_2=5$, công bội $q=2$. Số hạng thứ $7$  của cấp số nhân là
	 \choice
		{$u_7=320$} 
		{$u_7=640$} 
		{\True $u_7=160$} 
		 {$u_7=80$}
	\loigiai{ 
	Ta có $u_2=u_1q\Rightarrow  5=u_1\cdot 2\Rightarrow u_1=\dfrac{5}{2}$.\\
	 Vậy $u_7=u_1q^6=160$.
	} 
\end{ex}

% nguồn: Ok_TracNghiemLop11_Lan1
\begin{ex}%[1D2N3-4]%[ma:bbdcd7b1]
	Cho  cấp số nhân $(u_n)$ có $u_1=-\dfrac{1}{2}$ và $u_2=1$. Tính $u_{10}$.
	 \choice
		{$u_{10}=-256$} 
		 {\True $u_{10}=256$}
		{$u_{10}=-512$} 
		 {$u_{10}=512$}
	\loigiai{ 
		Cấp số nhân có công bội $q=\dfrac{u_2}{u_1}=-2$. Suy ra  $u_{10}=u_1q^9=-\dfrac{1}{2}(-2)^9=256.$
	 }
\end{ex}

% nguồn: Ok_TracNghiemLop11_Lan1
\begin{ex}%[1D2N3-4]%[ma:a6a74ca1]
	Cho cấp số nhân $(u_n)$ có $u_1=3$,  công bội $q=\dfrac{1}{3}$. Tìm $u_3$.
	 \choice
		{$\dfrac{1}{9}$} 
		 {$\dfrac{1}{27}$}
		 {$\dfrac{1}{6}$}
		 {\True $\dfrac{1}{3}$}
	\loigiai{Cấp số nhân $(u_n)$ có số hạng đầu $u_1$, công bội $q$ thì công thức số hạng tổng quát của cấp số  nhân là $u_n=u_1\cdot q^{n-1}$.\\
	Vậy $u_3=u_1\cdot  q^2=3\cdot\left(\dfrac{1}{3}\right)^2=\dfrac{1}{3}$.}
\end{ex}

% nguồn: Ok_TracNghiemLop11_Lan1
\begin{ex}%[1D2N3-4]%[ma:48a0a3b0]
	Cho cấp số nhân $\left(u_n\right)$ với $u_1=-2$ và $q=-5$. Viết  bốn số hạng đầu tiên của cấp số nhân.
	\choice 
		{$-2; 10; 50;-250$} 
		{\True $-2; 10;-50; 250$} 
		{$-2;-10;-50;-250$} 
		{$-2; 10;  50; 250$}
	 \loigiai{
		Ta có $\heva{& u_1=-2 \\& q=-5}\Leftrightarrow  \heva{&u_1=-2 \\
		 & u_2=u_1 q=10 \\
		& u_3=u_2 q=-50 \\ 
		&  u_4=u_3 q=250}$.
	 }
\end{ex}

% nguồn: Ok_Toan11_TracNghiem_Lan2
\begin{ex}%[1D2H3-4]%[ma:0a2c99a3]
	Xác định $x$ dương để $2x - 3$, $x$, $2x + 3$ lập thành cấp số nhân. 
	\choice 
		 { $x = 3$ }
		{\True $x = \sqrt{3}$  }
		{  $x = \pm \sqrt{3}$ }
		{ không có  giá trị nào của $x$ }
	\loigiai{ 
		Ta có $2x - 3$, $x$, $2x + 3$  lập thành cấp số nhân, nên ta có
		\[x^2 = (2x - 3)(2x + 3) \Leftrightarrow x^2 = 4x^2 - 9 \Leftrightarrow x^2 = 3 \Leftrightarrow x = \pm  \sqrt{3}.\]
		Vì $x$ dương nên $x =  \sqrt{3}$.
	 }
\end{ex}

% nguồn: Ok_Toan11_TracNghiem_Lan2
\begin{ex}%[1D2H3-4]%[ma:77b72fd4]
	Cho cấp  số nhân có $u_1=-3$, công bội $q=\dfrac{2}{3}$. Số $-\dfrac{96}{243}$ là số hạng thức mấy của dãy số này?
	 \choice
		 {Thứ $5$}
		{\True Thứ $6$} 
		{Thứ $7$} 
		 {Không phải là số hạng của dãy số này}
	 \loigiai{
	Giả sử $-\dfrac{96}{243}$ là số hạng thứ $n$ của dãy số này.\\ 
	Ta có  $u_1q^{n-1}=-\dfrac{96}{243}\Rightarrow (-3)\cdot\left(\dfrac{2}{3}\right)^{n-1}=-\dfrac{96}{243}\Rightarrow n=6$.\\
	Vậy $-\dfrac{96}{243}$  là số hạng thứ $6$ của dãy số này.
	} 
\end{ex}

% nguồn: Ok_Toan11_TracNghiem_Lan2
\begin{ex}%[1D2H3-4]%[ma:6a72d342]
	Cho cấp  số nhân $(u_n)$ có $u_1=1$, công bội $q=2$. Hỏi $2^{2\,022}$ là số hạng thứ mấy của $(u_n)$?
	 \choice
		{\True Số hạng thứ  $2\,023$}
		{Số  hạng thứ $2\,022$}
		{Số hạng thứ  $2\,021$}
		{Số hạng  thứ $2\,020$}
	\loigiai{ 
		Ta có $u_n=u_1q^{n-1}=1.2^{n-1}$.\\ 
		Khi đó $u_n=2^{2\,022}\Leftrightarrow n-1=2\,022\Leftrightarrow  n=2\,023$.\\
	Do đó $2^{2\,022}$ là số  hạng thứ $2\,023$ của $(u_n)$.}
\end{ex}

% nguồn: Ok_Toan11_TracNghiem_Lan2
\begin{ex}%[1D2H3-4]%[ma:c0ab69f0]
	Cho cấp số nhân $(u_n)$ có $u_1=1$, công bội $q=2$.  Số $1\,024$ là số hạng thứ mấy của dãy?
	 \choice
		 {\True $11$}
		 {$9$}
		{$8$} 
		{$10$} 
	 \loigiai{
	Ta có $u_n=u_1q^{n-1}\Leftrightarrow 2^{n-1}=1\,024\Leftrightarrow 2^{n-1}=2^{10}\Leftrightarrow  n-1=10\Leftrightarrow n=11$.
	 }
\end{ex}

% nguồn: Ok_Toan11_TracNghiem_Lan2
\begin{ex}%[1D2H3-4]%[ma:65cf5291]
	Cho tứ giác $ABCD$ có bốn góc tạo tành cấp số nhân có công bội $q=2$, góc có số đo nhỏ  nhất trong bốn góc đó là
	\choice 
		 {$1^{\circ}$}
		{$30^{\circ}$} 
		 {$12^{\circ}$}
		{\True  $24^{\circ}$}
	\loigiai{ 
		Giả sử bốn góc $A,\,B,\,C,\,D$  theo thứ tự lập thành cấp số nhân và $A$ nhỏ nhất.\\
		Khi đó $B=2A,\,C=4A,\,D=8A$.\\ 
	 Nên $A+2A+4A+8A=360^{\circ}\Rightarrow A=24^{\circ}$.}
\end{ex}

% nguồn: Ok_Toan11_TracNghiem_Lan2
\begin{ex}%[1D2H3-4]%[ma:5870d24c]
	Cho một cấp số nhân có số hạng thứ $4$ gấp $4\,096$ lần số hạng đầu tiên. Tổng hai số hạng  đầu của dãy số là $34$. Số hạng thứ $3$ của dãy số có giá trị bằng:
	 \choice
		{$1$} 
		{\True  $512$}
		{$1\,024$} 
		 {$32$}
	\loigiai{Theo bài ra ta  có
		\begin{eqnarray*}
			&\heva{&u_4 = 4\,096 \cdot  u_1\\&u_1+u_2=34} &\Leftrightarrow \heva{&u_1q^3 = 4\,096 \\&u_1(1+q)=34}\\
			& &\Leftrightarrow \heva{&q^3 = 4\,096 \\&u_1(1+q)=34}\\ 
			 & &\Leftrightarrow \heva{&q = 16 \\&u_1(1+q)=34}\\
			& &\Leftrightarrow  \heva{&q=16 \\&u_1=2.}
		\end{eqnarray*}
	Vậy $u_3=u_1q^2=2\cdot 16^2=512$.} 
\end{ex}

% nguồn: Ok_Toan11_TracNghiem_Lan2
\begin{ex}%[1D2H3-4]%[ma:34320333]
	Cho cấp số nhân $\left(u_n \right)$ có $u_1=1, q=2$.  Hỏi số $1\,024$ là số hạng thứ mấy?
	 \choice
		{\True  $11$}
		 {$9$}
		{$8$} 
		{$10$} 
	\loigiai{Ta  có $u_n=u_1.q^{n-1} \Leftrightarrow 1.2^{n-1}=1024 \Leftrightarrow 2^{n-1}=2^{10} \Leftrightarrow n-1=10 \Leftrightarrow n=11$.}
\end{ex}

% nguồn: Ok_Toan11_TracNghiem_Lan2
\begin{ex}%[1D2H3-4]%[ma:aaa16564]
	Cho cấp số nhân $(u_n)$ xác định bởi $\heva{&u_1=3\\ &u_{n+1}=-2u_n,\,\forall n  \in \mathbb{N}}$. Số $3\,072$ là số hạng thứ mấy?
	 \choice
		 {$9$}
		 {$10$}
		 {$12$}
		{\True  $11$}
	 \loigiai{
		Từ giả thiết $u_{n+1}=-2u_n,\,\forall n \in  {\mathbb{N}}^* \Rightarrow q=\dfrac{u_{n+1}}{u_n}=-2$.\\
		Áp dụng công thức tổng  quát của cấp số nhân $u_n=u_1\cdot q^{n-1},\forall n\geq 2$.\\
	Ta có  $3\,072=3\cdot \left(-2\right)^{n-1}\Leftrightarrow{\left(-2\right)^{n-1}}=\left(-2\right)^{10}\Leftrightarrow n-1=10\Leftrightarrow n=11$.}
\end{ex}

% nguồn: Ok_Toan11_TracNghiem_Lan2
\begin{ex}%[1D2H3-4]%[ma:29e62af9]
	Cho cấp số nhân $(u_n)$ có  $u_1=7,q=\dfrac{2}{3}$. Số $\dfrac{224}{243}$ là số hạng thứ mấy của $(u_n)$?
	\choice 
		 {$5$}
		{$7$} 
		 {$188$}
		 {\True $6$}
	\loigiai{ 
	Giả sử $u_n=\dfrac{224}{243}\Leftrightarrow 7\left(\dfrac{2}{3}\right)^{n-1}=\dfrac{224}{243}\Leftrightarrow{\left(\dfrac{2}{3}\right)^{n-1}}=\dfrac{32}{243}\Leftrightarrow{\left(\dfrac{2}{3}\right)^{n-1}}=\left(\dfrac{2}{3}\right)^5\Leftrightarrow n-1=5\Leftrightarrow  n=6$.}
\end{ex}

% nguồn: Ok_Toan11_TracNghiem_Lan2
\begin{ex}%[1D2H3-4]%[ma:78093c02]
	Cho cấp số nhân  $(u_n)$ có $u_1=3$ và $q=-2$. Số $-96$ là số hạng thứ mấy của cấp số nhân đã cho?
	 \choice
		{Số hạng thứ  $5$}
		{\True Số hạng  thứ $6$}
		{Số  hạng thứ $7$}
		{Không  là số hạng của cấp số đã cho}
	 \loigiai{
	Ta có $-96=u_n=u_1q^{n-1}=3\cdot (-2)^{n-1}\Leftrightarrow(-1)^{n-1} \cdot 2^{n-1}=-32=(-1)^5\cdot  2^5\Leftrightarrow n=6$.}
\end{ex}

% nguồn: Ok_Toan11_TracNghiem_Lan2
\begin{ex}%[1D2H3-4]%[ma:9863bc40]
	Cho cấp số nhân $(u_n)$, biết $u_1=1,\,q=2$. Hỏi  số $1\,024$ là số hạng thứ mấy?
	 \choice
		{\True Số hạng  thứ $11$}
		{Số  hạng thứ $9$}
		{Số  hạng thứ $8$}
		{Số hạng thứ  $10$}
	\loigiai{ 
	Ta có $1\,024=u_n=u_1\cdot  q^{n-1}=1 \cdot (2)^{n-1}\Leftrightarrow{2^{n-1}}=1\,024=2^{10}\Leftrightarrow n=11$.}
\end{ex}

% nguồn: Ok_Toan11_TracNghiem_Lan2
\begin{ex}%[1D2H3-4]%[ma:d2fdfdd7]
	Cho dãy $\left(u_n\right)$ với $u_n=\left(\dfrac{1}{2}\right)^n+1$, $\forall n\in{\mathbb{N}^*}$. Tính $S_{2\,019}=u_1+u_2+u_3+\cdots +u_{2\,019}$, ta  được kết quả
	 \choice
		 {\True $2\,020-\dfrac{1}{2^{2\,019}}$}
		 {$\dfrac{4\,039}{2}$}
		{$2\,019+\dfrac{1}{2^{2\,019}}$} 
		 {$\dfrac{6\,057}{2}$}
	 \loigiai{
		$S_{2\,019}=2\,019+\left(\dfrac{1}{2}\right)^1+\left(\dfrac{1}{2}\right)^2+\cdots +\left(\dfrac{1}{2}\right)^{2\,019}=2\,019+\dfrac{1}{2}\cdot  \dfrac{1-\left(\dfrac{1}{2}\right)^{2\,019}}{1-\dfrac{1}{2}}=2\,020-\dfrac{1}{2^{2\,019}}$.
	 }
\end{ex}

% nguồn: Ok_Toan11_TracNghiem_Lan2
\begin{ex}%[1D2H3-4]%[ma:e08665e2]
	Cho cấp số nhân $(u_n)$ với $u_1 = -2$ và $q = -5$. Viết bốn số hạng đầu tiên của  cấp số nhân.
	\choice 
		{$-2$;  $10$; $50$; $-250$}
		 {\True $-2$; $10$; $-50$; $250$}
		{$-2$; $-10$; $-50$; $-250$} 
		 {$-2$; $10$;$50$; $250$}
	 \loigiai{
		Vì $u_1 = -2$  và $q = -5$ nên ta có
		\begin{eqnarray*}
			u_1  &=& -2;\\
			u_2 &=& u_1 \cdot q = -2 \cdot  (-5) = 10;\\
			u_3 &=& u_2 \cdot q = 10 \cdot (-5)  = -50;\\
			u_4  &=& u_3 \cdot q = -50 \cdot (-5) = 250.
		\end{eqnarray*}
		Vậy dãy số là $-2$;  $10$; $-50$; $250$.
	 }
\end{ex}

% nguồn: Ok_Toan11_TracNghiem_Lan2
\begin{ex}%[1D2H3-4]%[ma:887f107e]
	Cho cấp số nhân $(u_n)$ có  $u_5=2$ và $u_9=6$. Tính $u_{21}$.
	\choice 
		{$18$} 
		{$54$} 
		{\True $162$} 
		 {$486$}
	\loigiai{ 
	Ta có $\heva{  & u_5=2 \\ &u_9=6 }\Rightarrow \heva{ & u_1q^4=2 \\ &u_1q^8=6 }\Rightarrow \heva{ & u_1=\dfrac{2}{3} \\ &q^4=3.}$\\
	Suy ra  $u_{21}=u_1q^{20}=u_1\left(q^4\right)^5=\dfrac{2}{3}\cdot 3^5=162$.
	} 
\end{ex}

% nguồn: Ok_Toan11_TracNghiem_Lan2
\begin{ex}%[1D2H3-4]%[ma:844d1c46]
	Cho cấp  số nhân $(u_n)$, $n \ge 1$ với công bội $q=2$ và có số hạng thứ hai $u_2=5$. Số hạng thứ $7$ của cấp số nhân là
	 \choice
		{$u_7=320$} 
		{$u_7=640$} 
		{\True $u_7=160$} 
		{$u_7=80$} 
	\loigiai{Ta có $(u_n)$, $n \ge 1$ là cấp số nhân có công bội  $q=2$ nên có số hạng tổng quát $u_n=q^{n-1}.u_1$.\\
	Vì $u_2=5=u_1.2 \Rightarrow  u_1=\dfrac{5}{2} \Rightarrow u_7=\dfrac{5}{2}.2^6=160$.\\ Vậy số hạng thứ $7$ của cấp số là $160$.}
\end{ex}

% nguồn: Ok_Toan11_TracNghiem_Lan2
\begin{ex}%[1D2H3-4]%[ma:ff7cb30a]
	 Cho dãy số $4$, $12$, $36$, $108$, $324$, $\ldots$. Số hạng thứ $10$ của dãy số đó là
	\choice 
		{$73\,872$} 
		{$77\,832$} 
		 {$72\,873$}
		 {\True $78\,732$}
	 \loigiai{
		Xét dãy số $4,12,36,108,324,\ldots$ là cấp số nhân có  $u_1=4$, $q=3$.\\
	Số hạng thứ $10$  của dãy số là $u_{10}=u_1\cdot q^9 =4 \cdot 3^9 =78\,732$.}
\end{ex}

% nguồn: Ok_Toan11_TracNghiem_Lan2
\begin{ex}%[1D2H3-4]%[ma:72b376b7]
	Cho cấp số nhân $\left(u_n \right)$ có số hạng đầu  $u_1=5$ và công bội $q=-2$. Số hạng thứ sáu của $\left(u_n \right)$ là
	\choice 
		 {$u_6=320$}
		 {\True $u_6=-160$}
		 {$u_6=-320$}
		{$u_6=160$} 
	\loigiai{Ta có $u_6=u_1.q^5 =  5.(-2)^5 = -160$.}
\end{ex}

% nguồn: Ok_Toan11_TracNghiem_Lan2
\begin{ex}%[1D2H3-4]%[ma:7df98bd4]
	Cho cấp số nhân  $(u_n)$ có công bội dương và $u_2=\dfrac{1}{4}$, $u_4=4$. Giá trị của $u_1$ là
	 \choice
		{$u_1=\dfrac{1}{6}$} 
		{\True  $u_1=\dfrac{1}{16}$}
		 {$u_1=-\dfrac{1}{16}$}
		 {$u_1=\dfrac{1}{2}$}
	 \loigiai{
	Ta có $\heva{ & u_2=\dfrac{1}{4} \\ &u_4=4  }\Rightarrow \heva{ & u_1q=\dfrac{1}{4} \\ &u_1q^3=4 }\Rightarrow q^2=16\Rightarrow q=4$ (vì $q>0$).\\
	Vậy  $u_1=\dfrac{u_2}{q}=\dfrac{1}{16}$.
	} 
\end{ex}

% nguồn: Ok_Toan11_TracNghiem_Lan2
\begin{ex}%[1D2H3-4]%[ma:8cf1ca27]
	Cho $\left(u_n \right)$ là cấp số nhân, công bội $q>0$. Biết $u_1=1,  u_3=4$. Tìm $u_4$.
	 \choice
		{$\dfrac{11}{2}$} 
		 {$2$}
		 {$16$}
		{\True  $8$}
	\loigiai{Ta có $\heva{&u_1=1\\&u_3=4\\&q>0} \Leftrightarrow \heva{&u_1=1\\&u_1 q^2=4\\&q>0}  \Leftrightarrow \heva{&u_1=1\\&q^2=4\\&q>0} \Leftrightarrow \heva{&u_1=1\\&q=2} \Rightarrow u_4=u_1q^3=8$.}
\end{ex}

% nguồn: Ok_TracNghiemLop11_Lan1
\begin{ex}%[1D2H3-4]%[ma:ad766a53]
	Cho cấp số nhân $\left(u_n\right)$ với  $u_1=-4$ và công bội $q=5$. Tính $u_4$.
	 \choice
		{$u_4=600$} 
		 {\True $u_4=-500$}
		 {$u_4=200$}
		{$u_4=800$} 
	\loigiai{ 
		Áp dụng công thức tính số hạng tổng quát  của cấp số nhân ta có
		 $$
		u_4=u_1 \cdot  q^3=(-4)\cdot 5^3=-500.
		$$ 
		 Vậy $u_4=-500$.
	} 
\end{ex}

% nguồn: Ok_TracNghiemLop11_Lan1
\begin{ex}%[1D2H3-4]%[ma:d54adfbd]
	Giả sử $\dfrac{\sin \alpha}{6},\, \cos \alpha,\, \tan \alpha$ theo thứ tự đó là một cấp  số nhân. Tính $\cos 2 \alpha$.
	 \choice
		{$-\dfrac{\sqrt{3}}{2}$} 
		 {$\dfrac{1}{2}$}
		{\True $-\dfrac{1}{2}$} 
		{$\dfrac{\sqrt{3}}{2}$} 
	 \loigiai{
		Điều kiện $\cos \alpha \neq 0 \Leftrightarrow \alpha \neq \dfrac{\pi}{2}+k \pi(k \in \mathbb{Z}$).\\ 
		 Theo tính chất của cấp số nhân, ta có
		$$\begin{aligned}
		&\cos ^2 \alpha=\dfrac{\sin  \alpha}{6} \cdot \tan \alpha\\
		&\Leftrightarrow 6 \cos ^2 \alpha=\dfrac{\sin ^2 \alpha}{\cos \alpha}\\ 
		&\Leftrightarrow 6 \cos ^3  \alpha-\sin ^2 \alpha=0\\
		&\Leftrightarrow 6 \cos ^3 \alpha+\cos  ^2 \alpha-1=0\\
		&\Leftrightarrow\cos  \alpha=\dfrac{1}{2}.
	\end{aligned}$$
	Ta có $\cos 2 \alpha=2 \cos ^2 \alpha-1=2  \cdot\left(\dfrac{1}{2}\right)^2-1=-\dfrac{1}{2}$.
} 
\end{ex}

% nguồn: Ok_TracNghiemLop11_Lan1
\begin{ex}%[1D2H3-4]%[ma:b403508e]
	Một cấp số nhân có số hạng thứ bảy bằng $\dfrac{1}{2}$, công  bội bằng $\dfrac{1}{4}$. Hỏi số hạng đầu tiên của cấp số nhân bằng bao nhiêu?
	 \choice
		 {$4096$}
		{\True  $2048$}
		 {$1024$}
		{$\dfrac{1}{512}$} 
	 \loigiai{
		Ta có $\heva{& q=\dfrac{1}{4}  \\ & \dfrac{1}{2}=u_7=u_1 q^6=\dfrac{u_1}{4^6}} \Rightarrow u_1=\dfrac{4^6}{2}=2048$.
	} 
\end{ex}

% nguồn: Ok_TracNghiemLop11_Lan1
\begin{ex}%[1D2H3-4]%[ma:506b8369]
	Cho cấp số  nhân $(u_n)$ có tổng $n$ số hạng đầu tiên là $S_n=5^n-1$. Tìm số hạng thứ 4 của cấp số nhân đã cho.
	\choice 
		{$u_4=100$} 
		 {$u_4=124$}
		 {\True $u_4=500$}
		{$u_4=624$} 
	 \loigiai{
		Ta có  $5^{n-1}-1=S_n=u_1 \cdot \dfrac{1-q^n}{1-q}=\dfrac{u_1}{q-1}\left(q^n-1\right) \Leftrightarrow \heva{&u_1=q-1 \\ &q=5 }\Leftrightarrow \heva{& u_1=4\\& q=5.}$\\
		Khi đó 
		 $u_4=u_1q^3=4\cdot 5^3=50$.
	 }
\end{ex}

% nguồn: Ok_TracNghiemLop11_Lan1
\begin{ex}%[1D2H3-4]%[ma:18261280]
	Cho cấp số nhân có $u_1=-8$, $q=\dfrac{3}{2}$. Số $-\dfrac{2187}{16}$ là số hạng thứ mấy của cấp  số này?
	 \choice
		{\True Thứ $8$} 
		{Không phải là số hạng  của cấp số}
		{Thứ  $9$}
		 {Thứ $7$}
	 \loigiai{
		 Giả sử $u_n=-\dfrac{2187}{32}$.\\
		Khi  đó $$u_n=u_1q^{n-1}\Rightarrow-\dfrac{2187}{16}=\left(-8\right)\cdot{\left(\dfrac{3}{2}\right)^{n-1}}\Rightarrow{\left(\dfrac{3}{2}\right)^{n-1}}=\dfrac{2187}{128}=\left(\dfrac{3}{2}\right)^7.$$
		Suy ra  $n-1=7\Rightarrow n=8$.\\
		Vậy $-\dfrac{2187}{16}$ là số hạng thứ  $8$.
	} 
\end{ex}

% nguồn: Ok_TracNghiemLop11_Lan1
\begin{ex}%[1D2H3-4]%[ma:924e98b6]
	Cho cấp số nhân $\left(u_n\right)$ có số hạng đầu $u_1=\dfrac{1}{2}$ và công  bội $q=3$. Tìm $u_5$
	\choice 
		 {\True $\dfrac{81}{2}$}
		{$\dfrac{163}{2}$} 
		{$\dfrac{27}{2}$} 
		 {$\dfrac{55}{2}$}
	 \loigiai{
		Áp dụng công thức số hạng tổng quát  $u_n=u_1 q^{n-1}(n \in \mathbb{N},\, n \geq 2)$.\\
		Do đó $u_5=u_1 q^4 \Rightarrow u_5=\dfrac{1}{2} \cdot 3^4=\dfrac{81}{2}$. 
	 }
\end{ex}

% nguồn: Ok_TracNghiemLop11_Lan1
\begin{ex}%[1D2H3-4]%[ma:2390c479]
	Cho cấp số nhân $(u_n)$ có số hạng đầu $u_3=8$ và $u_5=32$ và công bội $q>0$.  Số hạng thứ $10$ bằng bao nhiêu?
	\choice 
		{\True $1024$} 
		 {$\sqrt{33}$}
		{$512$} 
		{$-512$} 
	 \loigiai{
		Ta có $\heva{&u_3=8\\&u_5=32}\Rightarrow \heva{&u_1\cdot  q^2=8\\&u_1\cdot q^4=32}\Rightarrow q^2=4\Rightarrow \heva{& q=2 \text{ (do } q>0)\\&u_1=2.}$\\
		Suy  ra $u_{10}=u_1\cdot q^9=2\cdot 2^9=1\,024$.
	 }
\end{ex}

% nguồn: Ok_TracNghiemLop11_Lan1
\begin{ex}%[1D2H3-4]%[ma:ded96a10]
	Cho dãy số $(u_n)$ được xác định bởi $\heva{&u_1=1\\ &u_{n+1}=u_n+n, \forall n\geq 1}$. Tìm số các giá trị của $n$ để $u_n\leq 21\,529$. 
	\choice 
		{\True  $208$}
		 {$209$}
		{$210$} 
		{$211$} 
	 \loigiai{
		Chú ý rằng $u_{n+1}-u_n=n$, với mọi  $n\ge 1$ nên
		\begin{align*}
			 u_n&=\left(u_n-u_{n-1}\right)+ \left(u_{n-1}-u_{n-2}\right)+\cdots +(u_2-u_1)+u_1\\
			 &=(n-1)+(n-2)+\cdots +1+1\\
			&=\dfrac{(n-1)n}{2}+1=\dfrac{n^2-n+2}{2}. 
		\end{align*}
		Suy ra 
		\begin{align*}
			u_n\le 21\,529 \Leftrightarrow \dfrac{n^2-n+2}{2}\le  21\,529\Leftrightarrow n^2-n-43056\le 0\Leftrightarrow n\in [-207; 208].
		\end{align*}
		Do $n\in \mathbb{N^*}$ nên $n\in \{1; 2 3;\ldots;  208\}$.
	 }
\end{ex}

% nguồn: Ok_TracNghiemLop11_Lan1
\begin{ex}%[1D2H3-4]%[ma:cc8d0a7f]
	Cho cấp số nhân $(u_n)$ có tổng $n$ số hạng đầu  tiên là $S_n=\dfrac{3^n-1}{3^{n-1}}$. Tìm số hạng thứ 5 của cấp số nhân đã cho.
	 \choice
		{\True  $u_5=\dfrac{2}{3^4}$}
		{$u_5=\dfrac{1}{3^5}$} 
		 {$u_5=3^5$}
		 {$u_5=\dfrac{5}{3^5}$}
	 \loigiai{
		 Ta có $\dfrac{3^n-1}{3^{n-1}}=31-\dfrac{1}3^n=S_n=\dfrac{u_1}{1-q}\left(1-q^n\right) \Leftrightarrow \heva{&u_1=3(1-q) \\& q=\dfrac{1}{3}}\Leftrightarrow \heva{&u_1=2 \\& q=\dfrac{1}{3}}.$.\\
		Khi đó 
		$u_5=u_1q^4=\dfrac{2}{3^4}$. 
	} 
\end{ex}

% nguồn: Ok_Toan11_TracNghiem_Lan2
\begin{ex}%[1D2V3-4]%[ma:089522db]
	Một tứ giác lồi có số đo các góc lập thành một cấp số nhân. Biết rằng số đo của góc nhỏ nhất bằng $\dfrac{1}{9}$ số đo của góc nhỏ thứ ba. Hãy tính số đo của các góc trong tứ  giác đó.
	 \choice
		{$5^\circ$, $ 15^\circ$, $ 45^\circ$, $  225^\circ$}
		{\True $9^\circ$, $ 27^\circ$, $  81^\circ$, $ 243^\circ$}
		{$7^\circ$, $ 21^\circ$, $ 63^\circ$, $ 269^\circ$} 
		{$8^\circ$, $ 32^\circ$, $  72^\circ$, $ 248^\circ$}
	 \loigiai{
		Gọi số đo các góc của tứ giác là $a$, $aq$, $aq^2$, $aq^3$, trong đó $q >  1$.\\
		Theo giả thiết,  ta có $a=\dfrac{1}{9} aq^2$ nên $q=3$.\\
		Suy ra các góc của tứ giác là  $a$, $3a$, $9a$, $27a$.\\
		Vì tổng các góc trong tứ  giác bằng $360^\circ$ nên ta có $a+3a+9a+27a=360^\circ \Leftrightarrow a=9^\circ$.
	} 
\end{ex}

% nguồn: Ok_Toan11_TracNghiem_Lan2
\begin{ex}%[1D2V3-4]%[ma:9df77313]
	Cho một cấp số nhân có số hạng thứ $4$ gấp $4\,096$ lần số hạng đầu tiên. Tổng hai số hạng đầu tiên là $34$. Số hạng thứ $3$ của dãy số  có giá trị bằng
	\choice 
		 {$1$}
		 {\True $512$}
		{$1\,024$} 
		{$32$} 
	 \loigiai{
		Theo bài ra ta  có $\heva{&u_4=4\,096 \cdot u_1\\&u_1+u_2=34}\Leftrightarrow \heva{&q^3=4\,096\\&
		u_1\cdot (1+q)=34}\Leftrightarrow \heva{&q=16\\&17\cdot  u_1=34}\Leftrightarrow \heva{&q=16\\
		 & u_1=2.}$\\
	Vậy  $u_3=u_1\cdot q^2=2 c\cdot 16^2=512$.}
\end{ex}

% nguồn: Ok_Toan11_TracNghiem_Lan2
\begin{ex}%[1D2V3-4]%[ma:15982080]
	Cho cấp  số nhân $(u_n)$ với $u_1=-1,\,q=\dfrac{-1}{10}$. Số $\dfrac{1}{10^{103}}$ là số hạng thứ mấy của $(u_n)$?
	\choice 
		{Số hạng  thứ$103$}
		 {\True Số hạng thứ $104$}
		{Số hạng thứ $105$} 
		{Không là số hạng của cấp số đã  cho}
	\loigiai{ 
	Ta có $u_n=u_1\cdot q^{n-1}\Rightarrow\dfrac{1}{10^{103}}=-1  \cdot \left(-\dfrac{1}{10}\right)^{n-1}\Rightarrow n-1=103\Rightarrow n=104$.}
\end{ex}

% nguồn: Ok_Toan11_TracNghiem_Lan2
\begin{ex}%[1D2V3-4]%[ma:a78dec39]
	Cho cấp số nhân $(u_n)$ có tổng $n$ số hạng đầu tiên là $S_n=5^n-1$ với $n=1,2,\ldots$. Tìm số hạng đầu $u_1$ và công bội $q$ của cấp số nhân đó? 
	 \choice
		{$u_1=5$,  $q=4$}
		{$u_1=5$, $q=6$} 
		{\True  $u_1=4$, $q=5$}
		{$u_1=6$,  $q=5$}
	\loigiai{ 
	Ta có $\heva{&u_1=S_1=5-1=4\\&u_1+u_2=S_2=5^2-1=24}\Rightarrow \heva{&u_1=4\\&u_2=24-u_1=20}\Rightarrow u_1=4$,  $q=\dfrac{u_2}{u_1}=5$.}
\end{ex}

% nguồn: Ok_Toan11_TracNghiem_Lan2
\begin{ex}%[1D2V3-4]%[ma:0ea69ddc]
	Cho cấp số nhân $(u_n)$ thỏa mãn $\heva{&u_1-u_3+u_5=65\\&u_1+u_7=325}$. Tính  $u_3$.
	 \choice
		 {$u_3=15$}
		 {$u_3=25$}
		 {$u_3=10$}
		 {\True $u_3=20$}
	 \loigiai{
		Ta có: $\heva{&u_1-u_3+u_5=65\\ &u_1+u_7=325}  \Leftrightarrow\heva{&u_1-u_1\cdot q^2+u_1.q^4=65\\&u_1+u_1\cdot q^6=325}\Leftrightarrow\heva{&u_1(1-q^2+q^4)=65\quad (1)\\&u_1(1+q^6)=325.\quad (2)}$\\
		Chia từng  vế của $(1)$ cho $(2)$ ta được phương trình \\
		 $\dfrac{1-q^2+q^4}{1+q^6}=\dfrac{1}{5}\Leftrightarrow{q^6}-5q^4+5q^2-4=0.\quad (*)$\\
		Đặt $t=q^2,t\geq  0$.\\
		Phương trình $(*)$ trở  thành
		\[t^3-5t^2+5t-4=0 \Leftrightarrow(t-4)(t^2-t+1)=0 \Leftrightarrow  \hoac{&t=4\\ &t^2-t+1=0\,(\text{vô nghiệm}).}\]
		Với $t=4\Rightarrow q^2=4\Leftrightarrow q=\pm  2$.\\
		Với $q=\pm 2$ thay  vào $(2)$ ta được $u_1=5$.\\
	 Vậy $u_3=u_1 \cdot q^2=5 \cdot 4=20$.}
\end{ex}

% nguồn: Ok_Toan11_TracNghiem_Lan2
\begin{ex}%[1D2V3-4]%[ma:f97095be]
	Xen giữa số  $3$ và số $768$ là $7$ số để được một cấp số nhân có $u_1=3$. Khi đó $u_5$ là
	 \choice
		{$72$} 
		{$-48$} 
		 {$\pm 48$}
		 {\True $48$}
	 \loigiai{
		Ta có $u_1=3$ và $u_9=768$ nên $768=3\cdot q^8$  $\Rightarrow{q^8}=256$ $\Rightarrow q=\pm 2$.\\
	 Do đó $u_5=u_1\cdot q^4=3 \cdot 2^4=48$.}
\end{ex}

% nguồn: Ok_Toan11_TracNghiem_Lan2
\begin{ex}%[1D2V3-4]%[ma:5c01da00]
	Cho cấp số nhân $(u_n)$ với  $u_1=3,\,q=-2$. Số $192$ là số hạng thứ mấy của $(u_n)$?
	\choice 
		 {Số hạng thứ $5$}
		 {Số hạng thứ $6$}
		{\True Số hạng  thứ $7$}
		{Không là số hạng của cấp số đã  cho}
	 \loigiai{
	Ta có $u_n=u_1\cdot  q^{n-1}\Rightarrow 192=3 \cdot (-2)^{n-1}\Rightarrow(-2)^{n-1}=64\Rightarrow n-1=6\Rightarrow n=7$.}
\end{ex}

% nguồn: Ok_Toan11_TracNghiem_Lan2
\begin{ex}%[1D2V3-4]%[ma:d9e397a7]
	Cho ba số $x$; $5$; $2y$ theo thứ tự lập thành cấp số cộng và ba số $x$; $4$; $2y$ theo thứ tự lập thành cấp số nhân  thì $\left|x-2y\right|$ bằng
	\choice 
		{$\left|x-2y\right|=10$} 
		{$\left|x-2y\right|=9$} 
		 {\True $\left|x-2y\right|=6$}
		 {$\left|x-2y\right|=8$}
	 \loigiai{
		Do ba số $x$; $5$; $2y$ theo thứ  tự lập thành cấp số cộng, ta có $S=x+2y=10.\quad(1)$\\
		Ta lại có ba số $x$; $4$; $2y$ theo thứ tự lập thành cấp số nhân  nên $P=x\cdot 2y=16.\quad(2)$\\
		Từ $(1)$, $(2)$ suy ra  hai số $x$; $2y$ là nghiệm của phương trình $X^2-S\cdot X+P=0$ hay
		 $X^2-10X+16=0\Rightarrow \hoac{&X=2\\&X=8.}$
		Theo  yêu cầu bài toán $\left|x-2y\right|=\left|2-8\right|=6$.
	} 
\end{ex}

% nguồn: Ok_Toan11_TracNghiem_Lan2
\begin{ex}%[1D2V3-4]%[ma:d98fb5da]
	Cho cấp số nhân $(u_n)$  có $\heva{&u_2+u_7=198\\&u_3+u_8=396}$. Khi đó công bội của cấp số nhân $(u_n)$ là
	 \choice
		{\True  $2$}
		 {$3$}
		 {$4$}
		{$1$} 
	\loigiai{ 
		Gọi $q$ là  công bội của cấp số nhân $(u_n)$. Khi đó\\
		$\heva{&u_2+u_7=198\\ &u_3+u_8=396}\Leftrightarrow \heva{&u_1q+u_1q^6=198\\&u_1q^2+u_1q^7=396} 
	\Leftrightarrow \heva{&u_1q(1+q^5)=198\\& u_1q^2(1+q^5)=396}\Rightarrow  q=2$.}
\end{ex}

% nguồn: Ok_Toan11_TracNghiem_Lan2
\begin{ex}%[1D2V3-4]%[ma:62b7f020]
	Cho cấp số  nhân $(u_n)$ biết $\heva{&u_4-u_2=54\\&u_5-u_3=108}$. Tìm số hạng đầu $u_1$ và công bội $q$ của cấp số nhân trên.
	 \choice
		{\True $u_1=9$,  $q=2$}
		{$u_1=9$,  $q=-2$}
		 {$u_1=-9$, $q=-2$}
		 {$u_1=-9$, $q=2$}
	\loigiai{ 
		Ta có $\heva{&u_4-u_2=54\\&u_5-u_3=108} \Leftrightarrow\heva{&u_1\cdot q^3-u_1\cdot  q=54\\
		&u_1\cdot q^4-u_1  \cdot q^2=108}\Leftrightarrow\heva{&u_1 \cdot q(q^2-1)=54\\
		&u_1\cdot q^2(q^2-1)=108}\Leftrightarrow  \heva{&u_1=9\\&q=2.}$\\
	Vậy $u_1=9$,  $q=2$.}
\end{ex}

% nguồn: Ok_Toan11_TracNghiem_Lan2
\begin{ex}%[1D2V3-4]%[ma:f5131e8c]
	Giả sử $\dfrac{\sin \alpha}{6}$, $\cos \alpha$, $\tan \alpha$ theo thứ tự đó là một cấp số  nhân. Tính $\cos 2\alpha$.
	\choice 
		{  $\dfrac{\sqrt{3}}{2}$ }
		{ $-\dfrac{\sqrt{3}}{2}$ } 
		 { $\dfrac{1}{2}$ }
		{ \True $-\dfrac{1}{2}$ } 
	\loigiai{ 
		Điều kiện $\cos \alpha \ne 0 \Leftrightarrow \alpha \ne \dfrac{\pi}{2} + k\pi\quad  (k \in \mathbb{Z})$.\\
		Theo tính chất của cấp  số nhân, ta có
		\begin{eqnarray*}
			&&\cos^2 \alpha = \dfrac{\sin  \alpha}{6} \cdot \tan \alpha = \dfrac{\sin^2 \alpha}{6\cos \alpha}\\
			&\Rightarrow& 6 \cos^3 \alpha  - \sin^2 \alpha = 0\\ &\Leftrightarrow& 6 \cos^3 \alpha + \cos^2 \alpha - 1 = 0\\
			 &\Leftrightarrow& \cos \alpha = \dfrac{1}{2}.
		\end{eqnarray*}
		Ta có: $\cos 2\alpha = 2 \cos^2 \alpha - 1 = 2 \cdot \left( \dfrac{1}{2} \right)^2 - 1 = \dfrac{1}{2} - 1  = -\dfrac{1}{2}$.
	 }
\end{ex}

% nguồn: Ok_Toan11_TracNghiem_Lan2
\begin{ex}%[1D2V3-4]%[ma:ea669c77]
	Cho cấp số nhân có số hạng đầu $u_1=-2$, công bội $q=\dfrac{3}{4}$. Số  $-\dfrac{81}{128}$ là số hạng thứ mấy của cấp số này?
	\choice 
		 {\True $5$}
		{$4$} 
		 {$6$}
		 {$3$}
	\loigiai{ 
		 Áp dụng công thức cấp số nhân
		\begin{eqnarray*}
			 &u_n&=u_1q^{n-1}\\
			\Rightarrow &-\dfrac{81}{128}&=-2\cdot  \left(\dfrac{3}{4}\right)^{n-1}\\\Leftrightarrow&{\left(\dfrac{3}{4}\right)^4}&=\left(\dfrac{3}{4}\right)^{n-1}\\
			\Leftrightarrow&  n&=5.
		\end{eqnarray*}
	} 
\end{ex}

% nguồn: Ok_Toan11_TracNghiem_Lan2
\begin{ex}%[1D2V3-4]%[ma:10997a2d]
	Cấp  số nhân $(u_n)$ có $\heva{&u_{20}=8u_{17}\\& u_1+u_5=272}$. Tìm $u_1$, biết rằng $u_1\leq 100$.
	 \choice
		{\True  $u_1=16$}
		{$u_1=2$} 
		{$u_1=-16$} 
		 {$u_1=-2$}
	 \loigiai{
		Ta có\\ 
		$\heva{&u_{20}=8u_{17}\\&u_1+u_5=272}\Leftrightarrow\heva{ u_1\cdot q^{19}=8u_1 \cdot  q^{16}\\
		u_1+u_1\cdot q^4=272}\Leftrightarrow\heva{&u_1{q^{16}}(q^3-8)=0\quad (1)\\ &u_1(1+q^4)=272.\quad (2)}$\\ 
		Từ $(2)$  suy ra $u_1\neq 0$ do đó $(1)\Leftrightarrow\hoac{&q=0\\&q=2.}$\\
		Nếu $q=0$ thì $(2)\Leftrightarrow u_1=272$ không thoả  điều kiện $u_1\leq 100$.\\
	Nếu $q=2$ thì $(2)\Leftrightarrow u_1=16$  thoả điều kiện $u_1\leq 100$.}
\end{ex}

% nguồn: Ok_Toan11_TracNghiem_Lan2
\begin{ex}%[1D2V3-4]%[ma:ff153029]
	Bốn góc của một tứ giác tạo thành cấp số nhân và góc lớn nhất gấp $27$ lần góc nhỏ  nhất. Tổng của góc lớn nhất và góc bé nhất bằng
	\choice 
		{$56^{\circ}$} 
		{$102^{\circ}$} 
		{\True $252^{\circ}$} 
		{$168^{\circ}$} 
	\loigiai{ 
		Giả  sử $4$ góc $A$, $B$, $C$, $D$ theo thứ tự đó lập thành cấp số nhân thỏa yêu cầu với công bội $q$. Ta có
		\[\heva{&{A+B+C+D=360^\circ} \\&{D=27A}} \Leftrightarrow \heva{&{A\left(1+q+q^2+q^3\right)=360^\circ}  \\&{Aq^3=27A}}\Leftrightarrow \heva{&{q=3} \\&{A=9^\circ} \\&{D=Aq^3=243^\circ}} \Rightarrow A+D=252^\circ.\]
	} 
\end{ex}

% nguồn: Ok_Toan11_TracNghiem_Lan2
\begin{ex}%[1D2V3-4]%[ma:581cee86]
	Cho cấp số nhân $(u_n)$ biết $\heva{&u_4+u_6=540\\&u_1+u_3=20}$.  Công bội $q$ bằng
	 \choice
		 {$6$}
		{$27$} 
		{$2$} 
		{\True  $3$}
	\loigiai{ 
	Ta có: $\heva{&u_4+u_6=540\\&u_1+u_3=20}\Leftrightarrow\heva{&u_1q^3\cdot(1+q^2)=540\\&u_1(1+q^2)=20}\Leftrightarrow  \heva{&q^3=27\\&u_1(1+q^2)=20} \Leftrightarrow\heva{&q=3\\ &u_1=2.}$}
\end{ex}

% nguồn: Ok_Toan11_TracNghiem_Lan2
\begin{ex}%[1D2V3-4]%[ma:b259f731]
	Cho cấp số nhân $(u_n)$ có tổng  $n$ số hạng đầu tiên là $S_n=6^n-1$. Tìm số hạng thứ năm của cấp số nhân đã cho.
	\choice 
		 {$120\,005$}
		{$6\,840$} 
		 {$7\,775$}
		{\True  $6\,480$}
	\loigiai{ 
		Cấp số nhân $(u_n)$ có số hạng đầu $u_1$ và công bội $q$.\\ 
		Do  $S_n=6^n-1$ nên $q\neq 1$. Khi đó $S_n=\dfrac{u_1(1-q^n)}{1-q}=6^n-1$.\\
		Ta có $S_1=\dfrac{u_1(1-q)}{1-q}=6-1\Leftrightarrow  u_1=5$.\\
		$S_2=\dfrac{u_1(1-q^2)}{1-q}=6^2-1\Leftrightarrow q=6$.\\ 
	Vậy $u_5=u_1\cdot  q^4=5\cdot 6^4=6\,480$.}
\end{ex}

% nguồn: Ok_Toan11_TracNghiem_Lan2
\begin{ex}%[1D2V3-4]%[ma:725e74e5]
	 Cho cấp số nhân $u_1=-1$, $u_6=0,000\,01$. Khi đó $q$ và số hạng tổng quát là
	 \choice
		 {$q=\dfrac{1}{10}$, $u_n=\dfrac{-1}{10^{n-1}}$}
		 {$q=\dfrac{-1}{10}$, $u_n=-10^{n-1}$}
		{\True $q=\dfrac{-1}{10}$, $u_n=\dfrac{\left(-1\right)^n}{10^{n-1}}$} 
		 {$q=\dfrac{1}{10}$, $u_n=\dfrac{1}{10^{n-1}}$}
	\loigiai{ 
		Ta  có $u_6=u_1 \cdot q^5=0,000\,01 \Leftrightarrow{q^5}=\dfrac{-1}{10^5} \Leftrightarrow q=\dfrac{-1}{10}$.\\
	$\Rightarrow{u_n}=u_1\cdot q^{n-1}=-1  \cdot \left(\dfrac{-1}{10}\right)^{n-1}=\dfrac{\left(-1\right)^n}{10^{n-1}}$.}
\end{ex}

% nguồn: Ok_Toan11_TracNghiem_Lan2
\begin{ex}%[1D2V3-4]%[ma:60eb836f]
	Cho dãy số $\left(u_n\right)$ xác định bởi $u_1=1$; $u_{n+1}=\dfrac{3}{2}\left(u_n-\dfrac{n+4}{n^2+3n+2}\right), n\ge 1$. Giá trị của $u_{50}$ gần nhất với số nào  dưới đây?
	\choice 
		{$-312\,540\,600$} 
		{$-312\,540\,500$} 
		 {\True $-212\,540\,500$}
		 {$-212\,540\,600$}
	\loigiai{ 
		Ta có 
		\allowdisplaybreaks 
		\begin{eqnarray*}
			 &&u_{n+1}=\dfrac{3}{2}\left(u_n-\dfrac{n+4}{n^2+3n+2}\right)\\
			&\Leftrightarrow&u_{n+1}=\dfrac{3}{2}\left(u_n-\dfrac{3}{n+1}+\dfrac{2}{n+2}\right)\\ 
			 &\Leftrightarrow&u_{n+1}-\dfrac{3}{n+2}=\dfrac{3}{2}\left(u_n-\dfrac{3}{n+1}\right).\quad(1)
		\end{eqnarray*}
		Đặt $v_n=u_n-\dfrac{3}{n+1},n\ge 1$, ta có $v_1=u_1-\dfrac{3}{2}=-\dfrac{1}{2}$ và từ $(1)$  thu được $v_{n+1}=\dfrac{3}{2}{v_n}$.\\
		Suy ra  dãy số $\left(v_n\right)$ là một cấp số nhân với công bội $q=\dfrac{3}{2}$, ta có $v_n=v_1.\left(\dfrac{3}{2}\right)^{n-1}=\left(-\dfrac{1}{2}\right).\left(\dfrac{3}{2}\right)^{n-1}$.\\
		Từ đó ta  được $u_n=\left(-\dfrac{1}{2}\right)\cdot \left(\dfrac{3}{2}\right)^{n-1}+\dfrac{3}{n+1}\Rightarrow{u_{50}}\approx-212\,540\,500$.
	 }
\end{ex}

% nguồn: Ok_Toan11_TracNghiem_Lan2
\begin{ex}%[1D2V3-4]%[ma:a6be61a3]
	Tìm số  hạng đầu $u_1$ của cấp số nhân $(u_n)$ biết rằng $u_1+u_2+u_3=168$ và $u_4+u_5+u_6=21$.
	\choice 
		{$u_1=24$} 
		{$u_1=\dfrac{1\,334}{11}$} 
		{\True  $u_1=96$}
		{$u_1=\dfrac{217}{3}$} 
	 \loigiai{
	Ta có $\heva{ & u_1+u_2+u_3=168 \\  &u_4+u_5+u_6=21 }\Leftrightarrow \heva{ & u_1\left(1+q+q^2\right)=168 \\ &u_1q^3\left(1+q+q^2\right)=21 }\Rightarrow q^3=\dfrac{1}{8}\Rightarrow q=\dfrac{1}{2}$.\\
	Vậy $u_1=\dfrac{168}{1+q+q^2}=96$. 
	} 
\end{ex}

% nguồn: Ok_Toan11_TracNghiem_Lan2
\begin{ex}%[1D2V3-4]%[ma:b15c6fa6]
	Cho dãy số $\left(u_n\right)$ được xác định bởi $u_1=2$; $u_n=2u_{n-1}+3n-1$. Tìm số hạng thứ $2\,019$ của dãy  số.
	\choice 
		{\True  $u_{2\,019}=5\cdot 2^{2\,019}-6\,062$}
		{$u_{2\,019}=5\cdot 2^{2\,019}+6\,062$} 
		{$u_{2\,019}=5\cdot 2^{2\,020}-6\,062$} 
		{$u_{2\,019}=5\cdot  2^{2\,020}+6\,062$}
	\loigiai{ 
		Ta có $u_n=2u_{n-1}+3n-1 \Leftrightarrow {u_n}+3n+5=2\left[u_{n-1}+3\left(n-1\right)+5\right]$,  với $n\ge 2$; $n\in\mathbb{N}$.\\
		Đặt $v_n=u_n+3n+5$, ta  có $v_n=2v_{n-1}$ với $n\ge 2$; $n\in\mathbb{N}$.\\
		Như vậy, $\left(v_n\right)$ là cấp số nhân với công bội $q=2$ và  $v_1=10$, do đó $v_n=10\cdot 2^{n-1}=5\cdot 2^n$.\\
		Do đó  $u_n+3n+5=5\cdot 2^n$, hay $u_n=5\cdot 2^n-3n-5$ với $n\ge 2$; $n\in\mathbb{N}$.\\
		Nên $u_{2\,019}=5\cdot 2^{2\,019}-6\,062$. 
	} 
\end{ex}

% nguồn: Ok_Toan11_TracNghiem_Lan2
\begin{ex}%[1D2V3-4]%[ma:c0502ffa]
	Cho dãy số xác định  bởi $u_1=1$, $u_{n+1}=\dfrac{1}{3}\left(2u_n+\dfrac{n-1}{n^2+3n+2}\right)$; $n\in \mathbb{N}^{*}$. Khi đó $u_{2\,018}$ bằng
	 \choice
		{\True  $u_{2\,018}=\dfrac{2^{2016}}{3^{2\,017}}+\dfrac{1}{2\,019}$}
		 {$u_{2\,018}=\dfrac{2^{2\,018}}{3^{2\,017}}+\dfrac{1}{2\,019}$}
		{$u_{2\,018}=\dfrac{2^{2\,017}}{3^{2\,018}}+\dfrac{1}{2\,019}$} 
		{$u_{2\,018}=\dfrac{2^{2\,017}}{3^{2\,018}}+\dfrac{1}{2\,019}$} 
	 \loigiai{
		Ta có 
		\begin{eqnarray*}
			u_{n+1}=\dfrac{1}{3} \left(2u_n+\dfrac{n-1}{n^2+3n+2}  \right)&=&\dfrac{1}{3}\left(2u_n+\dfrac{3}{n+2}-\dfrac{2}{n+1} \right)\\
			&=&\dfrac{2}{3} u_n+\dfrac{1}{n+2}-\dfrac{2}{3}\cdot \dfrac{1}{n+1}\\ 
			 &\Leftrightarrow & u_{n+1}-\dfrac{1}{n+2}=\dfrac{2}{3} \left(u_n-\dfrac{1}{n+1}\right). \quad(1)
		\end{eqnarray*}
		Đặt $v_n=u_n-\dfrac{1}{n+1}$, từ $(1)$ ta suy ra $v_{n+1}=\dfrac{2}{3}  v_n$.\\
		Do đó $(v_n)$ là cấp số nhân với $v_1=u_1-\dfrac{1}{2}=\dfrac{1}{2}$, công bội $q=\dfrac{2}{3}$.\\ 
		Suy  ra
		\begin{eqnarray*}
			v_n=v_1\cdot q^{n-1}=\dfrac{1}{2}\left(\dfrac{2}{3} \right)^{n-1}&\Leftrightarrow&  u_n-\dfrac{1}{n+1}=\dfrac{1}{2}\cdot \left(\dfrac{2}{3} \right)^{n-1}\\
			&\Leftrightarrow  & u_n=\dfrac{1}{2}\cdot \left(\dfrac{2}{3}\right)^{n-1}+\dfrac{1}{n+1}
		\end{eqnarray*}
		Vậy  $u_{2\,018}=\dfrac{1}{2}\cdot \left(\dfrac{2}{3}\right)^{2\,017}+\dfrac{1}{2\,019}=\dfrac{2^{2016}}{3^{2\,017}}+\dfrac{1}{2\,019}$.
	 }
\end{ex}

% nguồn: Ok_Toan11_TracNghiem_Lan2
\begin{ex}%[1D2V3-4]%[ma:821e9bdf]
	Cho cấp số nhân $(u_n)$ biết $u_1+u_5=51,\,u_2+u_6=102$. Hỏi số $12\,288$ là số hạng thứ mấy của cấp số nhân$(u_n)$? 
	 \choice
		 {Số hạng thứ $10$}
		{Số hạng  thứ $11$}
		 {Số hạng thứ $12$}
		{\True Số hạng thứ $13$} 
	\loigiai{ 
		Gọi  q là công bội của cấp số nhân đã cho.\\
		Theo đề bài, ta có $\heva{&u_1+u_5=51\\&u_2+u_6=102}\Leftrightarrow\heva{ &u_1\left(1+q^4\right)=51\\ & u_1q\left(1+q^4\right)=102} \Rightarrow \heva{&u_1=3\\&q=2}\Rightarrow{u_n}=3\cdot 2^{n-1}$.\\ 
	Mặt khác $u_n=12\,288\Leftrightarrow{3\cdot 2^{n-1}}=12\,288\Leftrightarrow{2^{n-1}}=2^{12}\Leftrightarrow n=13$.} 
\end{ex}

% nguồn: Ok_Toan11_TracNghiem_Lan2
\begin{ex}%[1D2V3-4]%[ma:debe74a5]
	Cho dãy số  $\left(u_n\right)$ xác định bởi $\heva{&u_1=1\\&u_{n+1}=2u_n+5}$. Tìm số hạng thứ $2\,020$ của dãy.
	 \choice
		{\True $u_{2\,020}=3\cdot  2^{2\,020}-5$}
		{$u_{2\,020}=3\cdot  2^{2\,019}+5$}
		{$u_{2\,020}=3\cdot 2^{2\,019}-5$} 
		{$u_{2\,020}=3\cdot 2^{2\,020}+5$} 
	 \loigiai{
		Đặt $u_n=v_n-5\Rightarrow v_{n+1}-5=2\cdot (v_n-5)+5\Rightarrow  v_{n+1}=2v_n$.\\
		Có $u_1=1\Rightarrow v_1=6\Rightarrow u_n+5=6\cdot  2^{n-1}\Rightarrow u_n=6\cdot 2^{n-1}-5$.\\
	 Vậy $u_{2\,020}=6\cdot 2^{2\,019}-5=3\cdot 2^{2\,020}-5$}
\end{ex}

% nguồn: Ok_Toan11_TracNghiem_Lan2
\begin{ex}%[1D2V3-4]%[ma:b01b0aa1]
	Một cấp số nhân hữu hạn có công bội $q=2$, số hạng thứ bốn bằng $-24$ và số hạng cuối bằng $-1\,572\,864$. Hỏi cấp số nhân đó có bao nhiêu số  hạng?
	 \choice
		 {$18$}
		{$19$} 
		 {\True $20$}
		{$21$} 
	 \loigiai{
		Gọi cấp số  nhân đó là $u_1;u_2;u_3;\ldots;u_n$.\\
		 Ta có $u_4=-24\Leftrightarrow u_1\cdot q^3=-24\Leftrightarrow u_1\cdot 2^3=-24\Leftrightarrow u_1=-3$.\\
		$u_n=-1\,572\,864\Leftrightarrow u_1\cdot  q^{n-1}=-1\,572\,864\Leftrightarrow (-3)\cdot 2^{n-1}=-1\,572\,864\Rightarrow n-1=19\Rightarrow n=20$.\\
		Vậy cấp số  nhân có $20$ số hạng.
	} 
\end{ex}

% nguồn: Ok_Toan11_TracNghiem_Lan2
\begin{ex}%[1D2V3-4]%[ma:8eae139f]
	Cho cấp số nhân $u_n$ có $u_2=\dfrac{1}{4}$,  $u_5=16$. Tìm công bội $q$ và số hạng đầu $u_1$.
	\choice 
		{$q=\dfrac{1}{2}$, $u_1=\dfrac{1}{2}$} 
		{$q=-\dfrac{1}{2}$,  $u_1=-\dfrac{1}{2}$}
		{$q=-4$,  $u_1=-\dfrac{1}{16}$}
		{\True $q=4$,  $u_1=\dfrac{1}{16}$}
	\loigiai{ 
		 Ta có $\heva{& u_2=\dfrac{1}{4}\\& u_5=16} \Leftrightarrow \heva{& u_1\cdot q=\dfrac{1}{4}&(1)\\& u_1\cdot q^4=16. &(2)}$\\
	Chia hai vế của $(2)$ cho  $(1)$ ta được $q^3=64 \Leftrightarrow q=4 \Rightarrow u_1=\dfrac{1}{16}$.}
\end{ex}

% nguồn: Ok_Toan11_TracNghiem_Lan2
\begin{ex}%[1D2C3-4]%[ma:c13eb66b]
	Cho dãy số $(u_n)$ thỏa mãn $\heva{& 3u_1 =  \sqrt{3u_1 - u_2} + u_2 + 6 \\ & u_{n+1} = 2u_n,~\forall n \in \mathbb{N}^*}$. Tìm giá trị nhỏ nhất của $n$ để $u_n \ge 2^{2021}$.
	 \choice
		 {$2021$}
		 {$1012$}
		 {$2022$}
		{\True  $1011$}
	\loigiai{ 
		Ta có: $u_{n+1}  = 2u_n \Rightarrow \dfrac{u_{n+1}}{u_n} = 2,~\forall n \in \mathbb{N}^*$ nên dãy $(u_n)$ là cấp số nhân với công bội $q = 2$. \\
		$\Rightarrow u_2  = 2u_1$. \\
		 Mà
		\begin{eqnarray*}
			&&3u_1 = \sqrt{3u_1 - u_2} + u_2  + 6\\
			&\Leftrightarrow& 3u_1 - u_2 - \sqrt{3u_1 - u_2} - 6 = 0 \\ 
			&\Leftrightarrow& \left(\sqrt{3u_1 - u_2}\right)^2 - \sqrt{3u_1  - u_2} - 6 = 0 \\
			&\Leftrightarrow& \heva{& \sqrt{3u_1 - u_2} = 2~\text{(nhận)}\\ & \sqrt{3u_1 -  u_2} = -3~\text{(loại)}}\\
			&\Leftrightarrow & 3u_1 -  u_2 = 4.
		\end{eqnarray*}
		Từ đó ta có \[\heva{& u_2 = 2u_1 \\ & 3u_1 - u_2 = 4} \Rightarrow  u_1 = 4.\]
		$\Rightarrow (u_n)$ là cấp số nhân với công bội $q = 2, u_1 = 4$. Nên số hạng tổng quát  là \\
		$u_n = 2 \cdot 4^{n-1} = 2 \cdot  2^{2(n-1)} = 2^{2n-1},~\forall n \in \mathbb{N}^*$. \\
		$u_n \ge 2^{2021} \Leftrightarrow 2^{2n-1} \ge 2^{2021} \Leftrightarrow 2n-1 \ge 2021  \Leftrightarrow n \ge 1011$. \\
		Vậy giá trị nhỏ nhất thỏa mãn là  $1011$.
	} 
\end{ex}

% nguồn: Ok_Toan11_TracNghiem_Lan2
\begin{ex}%[1D2C3-4]%[ma:c2316de9]
	 Cho dãy số $(u_n)$ xác định bởi $\heva{ & u_1=1 \\ &u_{n+1}=2u_n+5 }$. Tính số hạng thứ $2\,018$ của dãy số trên.
	\choice 
		 {\True $u_{2\,018}=6\cdot 2^{2\,017}-5$}
		 {$u_{2\,018}=6\cdot 2^{2\,018}-5$}
		{$u_{2\,018}=6\cdot  2^{2\,017}+1$}
		 {$u_{2\,018}=6\cdot 2^{2\,018}+5$}
	\loigiai{ 
	Ta  có $u_{n+1}=2u_n+5\Leftrightarrow u_{n+1}+5=2(u_n +5)$.\\
	Đặt  $v_n=u_n+5$ ta có $v_{n+1}=2v_n$.\\
	Khi đó $(v_n)$  là cấp số nhân với $v_1=u_1+5=6,q=2$.\\
	Suy ra $v_n=6\cdot 2^{n-1}\Rightarrow u_n+5=6\cdot  2^{n-1}$.\\
	Vậy  $u_{2\,018}=6\cdot 2^{2\,017}-5$.
	 }
\end{ex}


%%%%%%%%%%%%%%%% ĐÚNG/SAI %%%%%%%%%%%%%%%%

% nguồn: Ok_dungsai_lop11_Lan1
\begin{ex}%[1D2N3-4]%[ma:703925d8]
 Cho cấp  số nhân $(u_n)$ có công bội $q$ và cấp số cộng $(v_n)$ có công sai $d$.
 \choiceTF 
  {\True  Với $u_1=-2$, công bội $q=3$. Số hạng $u_2$ là $-6$}
  {Với $u_1=-2$, công bội $q=\dfrac{1}{2}$. Số hạng thứ mười của cấp  số nhân đã cho bằng $\dfrac{1}{512}$}
  {\True Với $v_1=2$, công sai  $d=-3$. Số hạng thứ $11$ của cấp số cộng đã cho bằng $-28$}
  {\True Với $v_1=81$, công sai $d=-9$.  Tổng chín số hạng đầu tiên của cấp số cộng $(v_n)$ bằng $405$}
  \loigiai{
  \begin{itemchoice}
   \itemch \textbf{Đúng.}\\  Ta có $u_2=u_1\cdot q=-2\cdot 3=-6$.
   \itemch \textbf{Sai.}\\ Ta có $u_{10}=u_1\cdot q^9=(-2)\cdot \left(\dfrac{1}{2}\right)^9=-\dfrac{1}{256}$. 
   \itemch \textbf{Đúng.}\\ Ta có $v_{11}=v_1+10d=2+10\cdot(-3)=-28$. 
   \itemch \textbf{Đúng.}\\  Tổng chín số hạng đầu tiên của cấp số cộng đã cho
    \[S_9=\dfrac{9(2v_1+8d)}{2}=\dfrac{9[2\cdot 81+8\cdot (-9)]}{2}=405.\]
  \end{itemchoice}
 } 
\end{ex}

% nguồn: Ok_dungsai_lop11_Lan2
\begin{ex}%[1D2H3-4]%[ma:481a159b]
	Cho dãy số $\left(u_n\right)$ với $\heva{&u_1=2\\&u_{n+1}=3u_n}$, $\forall n \in \mathbb{N}^*$. Các mệnh  đề sau đúng hay sai?
	\choiceTF 
		{\True $\left(u_n\right)$  là cấp số nhân có cộng bội là $q=3$}
		{$u_n=2\cdot 3^n$} 
		 {$S_7=2\,189$}
		{\True $39\,366$ là số hạng thứ $10$ của cấp số  nhân}
	 \loigiai{
		\begin{itemchoice}
			 \itemch \textbf{Đúng} \\
			Ta có $u_{n+1}=3 u_n$ $\left(\forall n \in \mathbb{N}^*\right) \Rightarrow\left(u_n\right)$ là cấp số nhân có  cộng bội là $q=3$.
			 \itemch \textbf{Sai} \\
			Ta có $u_n=u_1 \cdot q^{n-1}=2 \cdot 3^{n-1}$. 
			\itemch  \textbf{Sai} \\
			Ta có $S_7=u_1  \dfrac{1-q^7}{1-q}=2 \cdot \dfrac{1-3^7}{1-3}=2\,186$.
			\itemch \textbf{Đúng} \\ 
			Ta có $u_n=2 \cdot  3^{n-1} \Leftrightarrow 2 \cdot 3^{n-1}=39\,366 \Leftrightarrow n=10$.
		\end{itemchoice}
	} 
\end{ex}

% nguồn: Ok_dungsai_lop11_Lan2
\begin{ex}%[1D2H3-4]%[ma:de82f874]
	Cho dãy số $\left(u_n\right)$ xác định bởi $u_1=1$ và $u_{n+1}=3 u_n+4$ với mọi $n \geq 1$.  Xét tính đúng sai của các mệnh đề sau
	\choiceTF 
		 {\True Số hạng $u_2=7$}
		{Số hạng $u_5=119$} 
		{Đặt $v_n=u_n+a$. Khi $a=-2$ thì $\left(v_n\right)$ lập thành  một cấp số nhân}
		{\True Số hạng tổng quát $u_n=3^n-2$ $(n\geq  1)$}
	\loigiai{ 
		\begin{itemchoice}
			 \itemch \textbf{Đúng} \\
			 Số hạng $u_2=3u_1+4=7$.
			\itemch \textbf{Sai} \\ 
			Ta tìm được $u_3=3u_2+4=25$, $u_4=3u_3+4=79$  nên $u_5=3u_4+4=241$.
			 \itemch \textbf{Sai} \\
			Ta  có $u_{n+1}=3u_n+4 \Leftrightarrow u_{n+1}+2=3\left(u_n+2\right)$.\\
			Đặt  $v_n=u_n+2$, ta có $v_{n+1}=3 v_n$.\\
			Do đó $\left(v_n\right)$ là một  cấp số nhân với $v_1=u_1+2=3$ và công bội $q=3$.\\
			\itemch  \textbf{Đúng} \\
			Số hạng tổng quát của dãy $\left(v_n\right)$ là  $v_n=v_1 \cdot q^{n-1}=3 \cdot 3^{n-1}=3^n \Rightarrow u_n=3^n-2$.
		\end{itemchoice}
	} 
\end{ex}

% nguồn: Ok_dungsai_lop11_Lan2
\begin{ex}%[1D2H3-4]%[ma:80cbf889]
	Cho cấp số nhân $\left(u_n\right)$,  biết $u_1=1$, $u_2 \cdot u_4=81$ và công bội $q>0$. Xét tính đúng sai của các mệnh đề sau
	 \choiceTF
		{Số hạng  $u_3=8$}
		{Số hạng $u_2=9$} 
		{\True Công bội  của cấp số nhân bằng $q=3$}
		{Số $59\,049$ là số  hạng thứ $10$ của cấp số nhân}
	\loigiai{ 
		Ta có $u_2=u_1q$, $u_4=u_1q^3 \Rightarrow u_2 \cdot u_4=81 \Leftrightarrow u_1^2 \cdot q^4=81  \Leftrightarrow q=3$.\\
		Do  đó $u_n=u_1 q^{n-1}=3^{n-1}$.
		\begin{itemchoice}
			\itemch \textbf{Sai}  \\
			Số  hạng $u_3=3^2=9$.
			 \itemch \textbf{Sai} \\
			 Số hạng $u_2=3^1=3$.
			\itemch  \textbf{Đúng} \\
			Công bội $q=3$. 
			 \itemch \textbf{Sai} \\
			Xét phương trình $u_n=59\,049  \Leftrightarrow 3^{n-1}=3^{10} \Leftrightarrow n=11$.
		\end{itemchoice}
	 }
\end{ex}

% nguồn: Ok_dungsai_lop11_Lan2
\begin{ex}%[1D2H3-4]%[ma:b0f77ddd]
	Cho tứ giác $ABCD$ có bốn góc tạo thành một cấp số nhân có  công bội bằng $2$. Khi đó
	\choiceTF 
		{\True Số đo góc nhỏ nhất bằng $24^\circ$} 
		{Số đo góc lớn nhất  bằng $196^\circ$}
		{Tổng số đo góc lớn nhất với góc nhỏ nhất bằng  $220^\circ$}
		{\True Số đo góc lớn nhất trừ cho số đo góc nhỏ nhất bằng $168^\circ$} 
	 \loigiai{
		Đặt $u_1$; $u_2$; $u_3$; $u_4$ theo thứ tự là số đo bốn  góc của tứ giác $ABCD$, gọi $q$ là công bội của cấp số nhân $u_1;u_2;u_3;u_4 \Rightarrow q=2$.\\
		Ta có 
		\begin{eqnarray*}
			\heva{& u_1+u_2+u_3+u_4=360^\circ \\ 
			&q=2} &\Leftrightarrow& \heva{&u_1\cdot  \dfrac{1-q^4}{1-q}=360^\circ \\
			&q=2} \\ 
			&\Leftrightarrow& \heva{&u_1\cdot  \dfrac{1-2^4}{1-2}=360^\circ \\
			 &q=2}\\
			 &\Leftrightarrow& \heva{&u_1=24^\circ\\
			 &q=2.}
		\end{eqnarray*}
		Vậy số đo bốn góc của  tứ giác $ABCD$ là $24^\circ;48^\circ;96^\circ;192^\circ$.
		\begin{itemchoice}
			\itemch \textbf{Đúng.}  \\
			 Số đo góc nhỏ nhất bằng $24^\circ$.
			\itemch \textbf{Sai.}  \\
			Số đo  góc lớn nhất bằng $192^\circ$.
			\itemch \textbf{Sai.} \\ 
			Tổng số đo góc lớn nhất với  góc nhỏ nhất bằng $216^\circ$.
			\itemch \textbf{Đúng.}  \\
			Số đo góc lớn nhất trừ  cho số đo góc nhỏ nhất bằng $168^\circ$.
		\end{itemchoice}
	} 
\end{ex}

% nguồn: Ok_dungsai_lop11_Lan2
\begin{ex}%[1D2H3-4]%[ma:a62676ca]
	Cho cấp số nhân $\left(u_n \right)$ với công bội $q < 0$ và $u_2=4,u_4=9$. Khi  đó
	 \choiceTF
		{\True Số hạng đầu $u_1=-\dfrac{8}{3}$} 
		{Số hạng $u_5=\dfrac{27}{2}$} 
		{$-\dfrac{2187}{32}$ là số hạng thứ $8$} 
		{\True Cấp số nhân có công  bội $q=-\dfrac{3}{2}$}
	\loigiai{ 
		Ta có $u_2=u_1 q=4,u_4=u_1 q^3=9\Rightarrow \dfrac{u_4}{u_2}=\dfrac{u_1 q^3}{u_1 q} \Rightarrow \dfrac{9}{4}=q^2\Rightarrow q=-\dfrac{3}{2}\,\,  (q < 0)$.
		\begin{itemchoice}
			\itemch \textbf{Đúng.}  \\
			Thay $q=-\dfrac{3}{2}$ vào  $u_2$, ta được $u_1 \left(-\dfrac{3}{2} \right)=4\Rightarrow u_1=-\dfrac{8}{3}$.
			\itemch \textbf{Sai.}\\ 
			Cấp số nhân đã cho  có số hạng đầu $u_1=-\dfrac{8}{3}$ và công bội $q=-\dfrac{3}{2}$.\\
			Khi  đó $u_n=-\dfrac{8}{3}\cdot\left(-\dfrac{3}{2} \right)^{n-1}$.
			Vậy $u_5=-\dfrac{27}{2}$. 
			\itemch  \textbf{Sai.}\\
			Ta có $-\dfrac{2187}{32} \ne-\dfrac{8}{3}.\left(-\dfrac{3}{2} \right)^7$ nên không phải là số  hạng thứ $8$.
			\itemch  \textbf{Đúng.} \\
			 Cấp số nhân có công bội $q=-\dfrac{3}{2}$.
		\end{itemchoice}
	 }
\end{ex}

% nguồn: Ok_dungsai_lop11_Lan2
\begin{ex}%[1D2H3-4]%[ma:21380431]
	Cho cấp số nhân $\left(u_n\right)$  có số hạng đầu $u_1=-2$ và công bội $q=\dfrac{1}{2}$. Xét tính đúng sai của các mệnh đề sau
	\choiceTF 
		 {\True Số hạng $u_2=-1$}
		{\True  Số hạng $u_3=-\dfrac{1}{2}$}
		{Số  hạng tổng quát $u_n=-2^{1-n}$}
		 {\True Tỉ số $\dfrac{u_3}{u_7}=16$}
	\loigiai{ 
		\begin{itemchoice}
			 \itemch \textbf{Đúng} \\
			Số hạng  $u_2=u_1 \cdot q=-2 \cdot \dfrac{1}{2}=-1$.
			 \itemch \textbf{Đúng} \\
			Số hạng $u_3=u_1  \cdot q^2=-2 \cdot\left(\dfrac{1}{2}\right)^2=-\dfrac{1}{2}$.
			\itemch \textbf{Sai} \\ 
			Số  hạng tổng quát $u_n=u_1 \cdot q^{n-1}=-2 \cdot\left(\dfrac{1}{2}\right)^{n-1}=-2^{2-n}$.
			\itemch  \textbf{Đúng} \\
			Tỉ  số $\dfrac{u_3}{u_7}=\dfrac{-\tfrac{1}{2}}{-\tfrac{1}{32}}=16$.
		\end{itemchoice}
	} 
\end{ex}

% nguồn: Ok_dungsai_lop11_Lan2
\begin{ex}%[1D2H3-4]%[ma:6c2913ed]
	Cho cấp số nhân $\left(u_n\right)$ có $u_4=40$ và $u_6=160$. Xét  tính đúng sai của các mệnh đề sau
	 \choiceTF
		 {Số hạng $u_5^2=4$}
		 {Số hạng $u_3=20$}
		{Công bội của cấp  số nhân bằng $q=-2$}
		{\True Số  hạng $u_{2n}>0$}
	\loigiai{ 
		Ta có $u_4=u_1q^3$, $u_6=u_1 q^5 \Rightarrow \dfrac{u_6}{u_4}=q^2 \Leftrightarrow  q^2=4 \Leftrightarrow q=\pm 2$.\\
		Với  $q=2$ thì $u_1=5$. Do đó $u_n=u_1q^{n-1}=5\cdot 2^{n-1}$.\\
		Với $q=-2$ thì  $u_1=-5$. Do đó $u_n=u_1q^{n-1}=-5 \cdot(-2)^{n-1}$.
		\begin{itemchoice}
			\itemch \textbf{Sai}  \\
			Số hạng  $u_5^2=u_4 \cdot u_6=2\,400$.
			\itemch  \textbf{Sai} \\
			Số hạng $u_3= \pm  5\cdot 4= \pm 20$.
			\itemch \textbf{Sai}  \\
			Công bội của cấp số nhân là $q=  \pm 2$.
			 \itemch \textbf{Đúng} \\
			Số hạng $u_{2n}=\dfrac{5}{2} \cdot  4^n>0$.
		\end{itemchoice}
	 }
\end{ex}

% nguồn: Ok_dungsai_lop11_Lan2
\begin{ex}%[1D2H3-4]%[ma:460f1097]
	Cho dãy số $\left(u_n\right)$ với $u_n=2\cdot 5^n$. Xét tính đúng sai của  các mệnh đề sau
	 \choiceTF
		{\True  Số hạng đầu $u_1=10$}
		{Số  hạng $u_5=160$}
		{$\left(u_n\right)$ là  cấp số cộng có công sai $d=5$ và số hạng đầu $u_1=2$}
		{\True $\left(u_n\right)$ là cấp số nhân  có công bội $q=5$ và số hạng đầu $u_1=10$}
	 \loigiai{
		\begin{itemchoice}
			\itemch \textbf{Đúng} \\ 
			 Ta có $u_1=2\cdot 5^1=10$.
			\itemch \textbf{Sai}  \\
			Ta có  $u_5=2\cdot 5^5=6\,250$.
			\itemch \textbf{Sai}  \\
			 Ta có $u_1=10$, $u_2=50$, $u_3=250$ nên $\left(u_n\right)$ không phải là cấp số cộng.
			\itemch  \textbf{Đúng} \\
			Ta có $\dfrac{u_{n+1}}{u_n}=\dfrac{2\cdot 5^{n+1}}{2\cdot 5^n}=5$ và $u_1=10$ nên $\left(u_n\right)$ là cấp  số nhân có công bội $q=5$ và số hạng đầu $u_1=10$.
		\end{itemchoice}
	 }
\end{ex}

% nguồn: Ok_dungsai_lop11_Lan2
\begin{ex}%[1D2H3-4]%[ma:b9378add]
	Cho cấp số nhân $\left(u_n \right)$ với công bội $q < 0$ và $u_2=4, u_4=9$. Xét tính đúng sai  của các khẳng định sau:
	 \choiceTF
		{\True Số hạng đầu  $u_1=-\dfrac{8}{3}$}
		{\True Cấp số nhân có công bội $q=-\dfrac{3}{2}$} 
		{Số hạng  $u_5=\dfrac{27}{2}$}
		{$-\dfrac{2187}{32}$ là số hạng thứ $8$} 
	\loigiai{ 
		\begin{itemchoice}
			\itemch  Ta có
			\[u_2 = u_1 q = 4, u_4 = u_1 q^3 = 9 \Rightarrow \dfrac{u_4}{u_2} = \dfrac{u_1 q^3}{u_1 q} = \dfrac{9}{4} \Rightarrow q^2 =  \dfrac{9}{4} \Rightarrow q = -\dfrac{3}{2}~(q < 0)\]
			Thay $q = -\dfrac{3}{2}$ vào $u_2$, ta được:$u_1 \left(-\dfrac{3}{2} \right) = 4 \Rightarrow u_1 =  -\dfrac{8}{3}$.
			\itemch Vậy cấp số nhân đã cho có số hạng  đầu $u_1 = -\dfrac{8}{3}$ và công bội $q = -\dfrac{3}{2}$.\\
			Khi  đó $u_n = -\dfrac{8}{3} \left(-\dfrac{3}{2} \right)^{n-1}$.
			\itemch Vậy $u_5  = -\dfrac{27}{2}$.
			\itemch $-\dfrac{2187}{32} \neq -\dfrac{8}{3} \left(-\dfrac{3}{2} \right)^7$ nên không phải là số hạng thứ $8$. 
		\end{itemchoice}
	} 
\end{ex}

% nguồn: Ok_dungsai_lop11_Lan2
\begin{ex}%[1D2V3-4]%[ma:7851dc91]
	Cho tứ giác $ABCD$ có bốn góc tạo thành một  cấp số nhân có công bội bằng $2$. Xét tính đúng sai của các khẳng định sau
	\choiceTF 
		{\True Số đo góc nhỏ nhất bằng  $24^\circ$}
		{Số đo góc lớn nhất bằng  $196^\circ$}
		{Tổng số đo  góc lớn nhất với góc nhỏ nhất bằng $220^\circ$}
		{\True Số đo góc lớn nhất trừ cho số đo góc nhỏ  nhất bằng $168^\circ$}
	 \loigiai{
		Đặt $u_1; u_2; u_3;  u_4$ theo thứ tự là số đo bốn góc của tứ giác $ABCD$, gọi $q$ là công bội của cấp số nhân, $u_1; u_2; u_3; u_4 \Rightarrow q = 2$.
		Ta  có
		\begin{eqnarray*}
			&& \heva{&u_1 +  u_2 + u_3 + u_4 = 360^\circ\\&q=2} \\
			&\Leftrightarrow& \heva{&u_1 \cdot \dfrac{1 - q^4}{1 - q} = 360^\circ\\&q=2} \\ 
			&\Leftrightarrow& \heva{& u_1 \cdot \dfrac{1 -  2^4}{1 - 2} = 360^\circ \\&q=2}\\
			 &\Leftrightarrow& \heva{&u_1=24^\circ\\&q=2.}
		\end{eqnarray*}
		Vậy số đo bốn góc của tứ giác  $ABCD$ là: $24^\circ$; $48^\circ$; $96^\circ$; $192^\circ$.
		\begin{itemchoice}
			\itemch 
			Số đo góc nhỏ nhất bằng  $24^\circ$.
			 \itemch
			Số đo góc  lớn nhất bằng $192^\circ$.
			 \itemch
			Tổng số đo góc lớn nhất  với góc nhỏ nhất bằng $24^\circ +192^\circ=216^\circ$.
			 \itemch
			Số  đo góc lớn nhất trừ cho số đo góc nhỏ nhất bằng $192^\circ-24^\circ=168^\circ$.
		\end{itemchoice}
	 }
\end{ex}

% nguồn: Ok_dungsai_lop11_Lan1
\begin{ex}%[1D2B3-4]%[ma:45e54885]
 Hãy cho biết tính đúng sai của mỗi mệnh đề  dưới đây?
  \choiceTF
  {\True Cấp số nhân $\left(u_n\right)$ có $u_1=3$, $q=2$. Khi đó $u_2=6$} 
   {Cho cấp số nhân $\left(u_n\right)$ với $u_1=3$, $q=\dfrac{1}{2}$. Khi đó $u_5=\dfrac{3}{32}$}
  {Dãy số $\left(u_n\right)$ là cấp số nhân có $10$ số hạng. Biết  số hạng đầu $u_1=7$ và công bội $q=-3$. Số hạng cuối của cấp số nhân bằng $-19683$}
  {\True Cấp số nhân  $(u_n)$ có $u_1=-3$, $\dfrac{u_8}{u_5}=125$. Khi đó $u_3=-75$}
 \loigiai{ 
  \begin{itemchoice}
   \itemch Ta  có $u_2=u_1q=6$.
   \itemch Ta có  $u_5=u_1q^4=\dfrac{3}{16}$.
   \itemch Ta có $u_{10}=u_1\cdot q^{10-1}=7\cdot  \left(-3\right)^9=-137781$.
    \itemch Ta có $\dfrac{u_8}{u_5}=125\Leftrightarrow \dfrac{u_1\cdot q^7}{ u_1\cdot q^4}=125\Leftrightarrow q^3=125\Leftrightarrow q=5$.\\
   Vậy  $u_3=u_1\cdot q^2=-3\cdot 5^2=-75$.
  \end{itemchoice}
 } 
\end{ex}

% nguồn: Ok_dungsai_lop11_Lan1
\begin{ex}%[1D2B3-4]%[ma:94c0c876]
 Hãy  cho biết tính đúng sai của mỗi mệnh đề dưới đây?
  \choiceTF
  {\True Cho cấp số nhân $\left(u_n\right)$ có  $u_1=3$, $q=-2$. Số $192$ là số hạng thứ bảy của cấp số nhân đã cho}
  {Cho cấp  số nhân $\left(u_{n}\right)$ có $u_1=1$, $q=2$. Số $256$ là số hạng thứ tám của cấp số nhân đã cho}
  {\True Cho cấp số nhân có $u_1=-1$, $q=-\dfrac{1}{10}$. Số $\dfrac{1}{10^{51}}$ là  số hạng thứ $52$}
  {Cho cấp  số nhân $\left(u_{n}\right):u_{1}=1$, $q=2$ . Số $2048$ là số hạng thứ $11$ của cấp số nhân đã cho}
 \loigiai{ 
  \begin{itemchoice}
   \itemch Số hạng thứ $n$ của cấp số nhân là $u_n=u_1\cdot q^{n-1}\Leftrightarrow  3\cdot (-2)^{n-1}=192\Rightarrow n=7$.\\
   Vậy số $192$ đúng là số hạng thứ $7$  của cấp số nhân đã cho.
   \itemch Giả  sử $u_n=256$. Khi đó $$u_1q^{n-1}=256\Leftrightarrow 1\cdot 2^{n-1}=2^8\Leftrightarrow n-1=8\Leftrightarrow n=9.$$
   Vậy $256$  là số hạng thứ $9$.
    \itemch Ta có $u_n=u_1\cdot q^{n-1}\Leftrightarrow \dfrac{1}{10^{51}}=-1\left(-\dfrac{1}{10}\right)^{n-1}\Leftrightarrow -1\left(-\dfrac{1}{10}\right)^{51}=-1\left(-\dfrac{1}{10}\right)^{n-1}\Leftrightarrow n=52$.
   \itemch Giả sử  $2048$ là số hạng thứ $n$ ta có $u_n=u_1\cdot q^{n-1}=1\cdot2^{n-1}=2048 \Leftrightarrow n-1=11 \Leftrightarrow n=12$.
  \end{itemchoice}
 } 
\end{ex}

% nguồn: Ok_dungsai_lop11_Lan1
\begin{ex}%[1D2B3-4]%[ma:2bc677ce]
 Hãy cho biết tính  đúng sai của mỗi mệnh đề dưới đây?
 \choiceTF 
  {Cho cấp số  nhân $(u_n)$ có $u_1=1$ và công bội $q=2$. Số hạng thứ $11$ của $(u_n)$ bằng $2048$}
  {\True Cho cấp số nhân $\left(u_n\right)$, biết $u_1=1$ và $u_2=-2$.  Khi đó $u_{10}=-512$}
  {\True Cho cấp số nhân  $(u_n)$ biết $u_1=-3$; $u_2=6$. Khi đó $u_5=-48$}
  {Một cấp số nhân có số hạng đầu tiên bằng  $2$ và số hạng thứ tư bằng $54$ thì số hạng thứ $6$ bằng $162$}
 \loigiai{ 
  \begin{itemchoice}
   \itemch Số hạng thứ $11$ của $(u_n)$ là $u_{11}=u_1\cdot  q^{10}=1\cdot 2^{10}=1024$.
   \itemch Vì $\left(u_n\right)$ là cấp số nhân nên $q=\dfrac{u_2}{u_1}=-2$.\\ 
   Vậy $u_{10}=u_1 \cdot q^9=1 \cdot  (-2)^9=-512$.
   \itemch Ta có công bội  $q=\dfrac{u_2}{u_1}=\dfrac{6}{-3}=-2$.\\
   Suy ra  $u_5=u_1\cdot q^4=-3\cdot (-2)^4=-48$.
   \itemch Ta có $ u_4=q^3u_1\Rightarrow q^3=\dfrac{u_4}{u_1}=\dfrac{54}{2}=27\Rightarrow q=3\Rightarrow u_6=q^5u_1=3^5\cdot  2=486 $.
  \end{itemchoice}
  }
\end{ex}


%%%%%%%%%%%%%%%% TRẢ LỜI NGẮN %%%%%%%%%%%%%%%%

% nguồn: Ok_TraLoiNgan_Lop11_Lan1
\begin{ex}%[1D2H3-4]%[ma:15f3678f]
	Khi nghiên cứu về một loại virus, người ta nhận thấy cứ sau mỗi phút, số  lượng virus tăng lên gấp ba lần trước đó. Giả sử ban đầu có $5$ con virus, gọi $n$ là số lượng virus có sau $11$ phút. Tính tổng tất cả các chữ số của $n$.
	\shortans{36} 
	\loigiai{ 
		Sau mỗi phút, số lượng virus tăng lên gấp $3$  lần trước đó.\\
		Suy ra số lượng con  vius có sau $11$ phút sẽ tăng thêm $3^{11}$ (lần).\\
		Sau $11$ phút, số lượng con virus là $5\cdot3^{11}=885\,735$  (con).\\
		Vậy tổng tất cả các chữ số của $n=885\,735$ là $8+8+5+7+3+5=36$. 
	 }
\end{ex}

% nguồn: Ok_TraLoiNgan_Lop11_Lan1
\begin{ex}%[1D2H3-4]%[ma:f4498c58]
	Cho cấp số nhân $(u_n)$, biết: $\heva{&u_2=6 \\& S_3=43.}$  Tích của các số hạng đầu và công bội bằng
	\shortans{36} 
	 \loigiai{
		Áp dụng công thức ${u_n=u_1 \cdot q^{n-1}}$ và  ${S_n=\dfrac{u_1 \cdot\left(1-q^n\right)}{1-q}}$.\\
		Ta  có
		$$\begin{aligned}
		& \heva{&u_2=6 \\&S_3=43}\Leftrightarrow \heva{&u_1q=6 \\&\dfrac{u_1\cdot (1-q^3)}{1-q}=43}\Leftrightarrow \heva{&u_1q=6 \\&u_1(1+q+q^2)=43}\Leftrightarrow \heva{&u_1=\dfrac{6}{q} \\ &\dfrac{q}{1+q+q^2}=\dfrac{6}{43}}\\ 
		 & \Rightarrow 43q=6.\left(1+q+q^2\right)\Leftrightarrow 6q^2-37q+6=0\Leftrightarrow \hoac{&q=6 \\&q=\dfrac{1}{6}.}
	\end{aligned}$$
	Với $q=6$ thì $u_1=\dfrac{6}{6}=1$.\\ 
	 Với $q=\dfrac{1}{6}$ thì $u_1=\dfrac{6}{\dfrac{1}{6}}=36$.\\
	Suy ra $\heva{&q=6 \\ &u_1=1}$ hoặc $\heva{&q=\dfrac{1}{6} \\  &u_1=36.}$\\
	Vậy tích các số  hạng đầu và công bội bằng $6\cdot 1\cdot \dfrac{1}{6}\cdot 36=36.$
} 
\end{ex}

% nguồn: Ok_TraLoiNgan_Lop11_Lan1
\begin{ex}%[1D2H3-4]%[ma:0281af73]
	Tìm $u_3$ của cấp  số nhân $\left(u_n\right)$, biết $\heva{&u_1+u_5=51 \\ &u_2+u_6=102}$.
	 \shortans{12}
	\loigiai{ 
		Áp dụng công thức  $u_n=u_1 \cdot q^{n-1}$.\\
		Ta có $\heva{&u_1+u_5=51 \\ &u_2+u_6=102} \Leftrightarrow\heva{u_1+u_1 q^4=51 \\ u_1 q+u_1 q^5=102} \Leftrightarrow\heva{u_1 \cdot\left(1+q^4\right)=51 \\ u_1  q \cdot\left(1+q^4\right)=102} \Leftrightarrow\heva{u_1=\dfrac{51}{1+q^4} \\ \frac{u_1 \cdot\left(1+q^4\right)}{u_1 q \cdot\left(1+q^4\right)}=\dfrac{51}{102}}$.\\
		Suy ra $\dfrac{1}{q}=\dfrac{1}{2} \Leftrightarrow  q=2$.\\
		Khi đó $u_1=\dfrac{51}{1+2^4}=3$.\\ 
		Vậy $u_1=3;\, q=2\Rightarrow u_3=u_1q^2=3\cdot  2^2=12$.
	} 
\end{ex}

% nguồn: Ok_TraLoiNgan_Lop11_Lan1
\begin{ex}%[1D2H3-4]%[ma:8952fec6]
	Cho cấp số nhân ${\left(u_n\right)}$ có tổng ${n}$ số hạng đầu tiên là ${S_n=5^n-1}$. Gọi số hạng đầu ${u_1}$  và công bội ${q}$ của cấp số nhân đó. Tính $u_2$.
	 \shortans{20}
	 \loigiai{
		Ta có ${u_1=S_1=5-1=4}$  và ${u_2=S_2-S_1=\left(5^2-1\right)-(5-1)=20}$.\\
		Vậy ${u_1=4,\, q=5}\Rightarrow u_2=u_1q=4\cdot 5=20$. 
	 }
\end{ex}

% nguồn: Ok_TraLoiNgan_Lop11_Lan1
\begin{ex}%[1D2H3-4]%[ma:6acd5f32]
	Cho cấp  số nhân $(u_n)$, biết: $\heva{&u_5+u_2=36 \\ &u_6-u_4=48.}$. Tính tổng số hạng đầu và công bội của cấp số nhân $(u_n)$.
	 \shortans{4}
	 \loigiai{
		Ta có: $\heva{&u_5+u_2=36 \\ &u_6-u_4=48}\Leftrightarrow\heva{&u_1 \cdot q^4+u_1 \cdot q=36 \\ &u_1 \cdot q^5-u_1 \cdot q^3=48}\Leftrightarrow\heva{&u_1\left(q^4+q\right)=36 \\ &u_1 \cdot\left(q^5-q^3\right)=48}\Leftrightarrow\heva{&u_1=\dfrac{36}{q^4+q} \\&  \dfrac{q^4+q}{q^5-q^3}=\dfrac{36}{48}}$
		$\begin{aligned}
		& \Rightarrow 4\cdot \left(q^4+q\right)=3\cdot  \left(q^5-q^3\right) \\
		& \Leftrightarrow  4q^4+4q=3q^5-3q^3\Leftrightarrow 4q^4+4q-3q^5+3q^3=0 \\
		&  \Leftrightarrow q\left(-3q^4+4q^3+3q^2+4q\right)=0\Leftrightarrow -q(q+1)(q-2) \\
		& \Leftrightarrow \hoac{&q=0 \\&q=-1 \\&q=2  \\&3q^2-q+2=0 \\}\\
		&  \Leftrightarrow q=2.\, (\text{do } q\ne 0;q\ne -1) \\
	\end{aligned}$\\
	 Với $q=2$ thì $u_1=\dfrac{36}{2^4+2}=2$
	$$\Leftrightarrow q\left(-3 q^4+4 q^3+3 q^2+4 q\right)=0 \Leftrightarrow-q(q+1)(q-2)\left(3 q^2-q+2\right)=0.$$ 
	Vậy  $q=2; u_1=2$, suy ra $u_1+q=2+2=4$.
} 
\end{ex}

% nguồn: Ok_TraLoiNgan_Lop11_Lan2
\begin{ex}%[1D2H3-4]%[ma:c289eed0]
	Tìm số hạng đầu của một cấp số nhân, biết rằng công bội là $3$, tổng  số các số hạng là $728$ và số hạng cuối là $486$.
	\shortans{2} 
	 \loigiai{
		Ta có 
		\begin{eqnarray*}
			&&\heva{&q=3\\&S_n=728\\&u_n=486} \Leftrightarrow \heva{& 
			q=3 \\& 
			u_1  \dfrac{1-q^n}{1-q}=728 \\&
			u_1 \cdot q^{n-1}=486} 
			\Leftrightarrow\heva{& 
			q=3 \\& u_1\left(1-q^n\right)=728(1-q)  \\&
			 u_1 q^n=486 q} \\
			 &\Leftrightarrow&\heva{&
			 q=3 \\& u_1-u_1 \cdot q^n=728(1-q)} \Leftrightarrow \heva{&u_1-486 q=728(1-q) \\&
			u_1 \cdot q^n=486 q}  \\
			&\Rightarrow& u_1-486\cdot  3=728(1-3) \Rightarrow u_1=2.
		\end{eqnarray*}
	} 
\end{ex}

% nguồn: Ok_TraLoiNgan_Lop11_Lan2
\begin{ex}%[1D2H3-4]%[ma:c5a476aa]
	Cho cấp số nhân $\left(u_n\right)$ biết $u_1+u_5=51$, $ u_2+u_6=102$. Hỏi số $12\,288$ là  số hạng thứ mấy của cấp số nhân $\left(u_n\right)$?
	 \shortans{13}
	 \loigiai{
		Gọi $q$ là công bội của cấp số nhân đã  cho. Theo đề bài, ta có
		\[\heva{&u_1+u_s=51\\&u_2+u_6=102} 
		 \Leftrightarrow \heva{&
		 u_1\left(1+q^4\right)=51 \\&
		 u_1 q\left(1+q^4\right)=102
		} \Rightarrow q=2 \Rightarrow  u_1=3 \Rightarrow u_n=3\cdot 2^{n-1}.\]
		Mặt khác $u_n=12\,288 \Leftrightarrow 3\cdot  2^{n-1}=12\,288 \Leftrightarrow 2^{n-1}=2^{12} \Leftrightarrow n=13$.
	 }
\end{ex}

% nguồn: Ok_TraLoiNgan_Lop11_Lan2
\begin{ex}%[1D2H3-4]%[ma:b8ac8c29]
	Cho cấp số nhân  $(u_n)$ có $u_1 = 3$ và $q = -2$. Số $192$ là số hạng thứ mấy của cấp số nhân đã cho?
	 \shortans{7}
	 \loigiai{
		$192 = u_n = u_1 q^{n-1} = 3  \cdot (-2)^{n-1} \Leftrightarrow (-1)^{n-1} \cdot 2^{n-1} = 64 = (-1)^6 \cdot 2^6 \Leftrightarrow n = 7$.
	 }
\end{ex}

% nguồn: Ok_TraLoiNgan_Lop11_Lan2
\begin{ex}%[1D2H3-4]%[ma:7d3f62da]
	Cho cấp số nhân có $u_1=\dfrac{1}{2}, q=-2$. Số $-1\,024$ là số hạng thứ mấy của cấp số nhân đã cho? 
	\shortans{12} 
	 \loigiai{
		Ta có 
		$u_n=-1\,024 \Leftrightarrow u_1 \cdot q^{n-1}=-1\,024 \Leftrightarrow \dfrac{1}{2} \cdot(-2)^{n-1}=-1\,024 \Leftrightarrow(-2)^{n-1}=(-2)^{11} \Leftrightarrow  n=12$.\\
		Vậy số $-1\,024$ là số hạng thứ $12$ của cấp số  nhân đã cho.
	} 
\end{ex}

% nguồn: Ok_TraLoiNgan_Lop11_Lan2
\begin{ex}%[1D2H3-4]%[ma:a928425a]
	Một cấp số nhân có hai số hạng liên tiếp là $16$ và $36$. Số  hạng tiếp theo bằng bao nhiêu?
	 \shortans{81}
	 \loigiai{
		Ta có  cấp số nhân $(u_n)$ có
		\[\heva{&u_k=16\\&u_{k+1}=36} \Rightarrow q=\dfrac{u_{k+1}}{u_k}=\dfrac{9}{4}\Rightarrow  u_{k+2}=u_{k+1}q=81.\]
	 }
\end{ex}

% nguồn: Ok_TraLoiNgan_Lop11_Lan2
\begin{ex}%[1D2H3-4]%[ma:27c482cd]
	Cho cấp số nhân $\left(u_n\right)$ có $u_1=-3$, công  bội $q=-2$. Hỏi $-192$ là số hạng thứ bao nhiêu cùa $\left(u_n\right)$?
	\shortans{7} 
	\loigiai{ 
		Ta có $u_n=u_1 \cdot q^{n-1} \Rightarrow-192=-3 \cdot(-2)^{n-1}  \Rightarrow(-2)^{n-1}=64 \Rightarrow n=7$.
	} 
\end{ex}

% nguồn: Ok_TraLoiNgan_Lop11_Lan2
\begin{ex}%[1D2H3-4]%[ma:580ecbd6]
	Cho cấp số nhân $\left(u_n\right)$ thỏa  mãn $\heva{&u_4=\dfrac{2}{27} \\& u_3=243 u_8}$. Số $\dfrac{2}{6\,561}$ là số hạng thứ bao nhiêu của cấp số nhân trên?
	 \shortans{9}
	 \loigiai{
		Gọi $q$ là  công bội của cấp số.\\
		Theo giả thiết ta có 
		$ 
		 \heva{&u_1q^3= \dfrac{2}{27}\\&u_1q^2=243\cdot u_1q^7}\Leftrightarrow \heva{&u_1q^3 =\dfrac{2}{27}\\&q^5=\dfrac{1}{243}} \Leftrightarrow \heva{&
		q=\dfrac{1}{3}  \\&u_1=2.}
		$\\ 
		Ta có  $u_n=\dfrac{2}{3^{n-1}} \Rightarrow u_n=\dfrac{2}{6\,561} \Leftrightarrow 3^{n-1}=6\,561=3^8 \Rightarrow n=9$.\\
		Vậy $\dfrac{2}{6\,561}$  là số hạng thứ $9$ của cấp số.
	 }
\end{ex}

% nguồn: Ok_TraLoiNgan_Lop11_Lan2
\begin{ex}%[1D2H3-4]%[ma:7b31327b]
	Cho cấp số nhân $(u_n)$ có $u_1 = -1$ và $q = -\dfrac{1}{10}$. Số $\dfrac{1}{10^{103}}$ là số hạng  thứ mấy của cấp số nhân đã cho?
	\shortans{104} 
	 \loigiai{
		\[ 
		\dfrac{1}{10^{103}} = u_n = u_1 q^{n-1}  = -1 \cdot \left(-\dfrac{1}{10}\right)^{n-1} = \dfrac{(-1)^n}{10^{n-1}} \Rightarrow
		\heva{&n \text{ chẵn}\\&n-1 =  103} \Leftrightarrow n = 104.
		\] 
	 }
\end{ex}

% nguồn: Ok_TraLoiNgan_Lop11_Lan2
\begin{ex}%[1D2H3-4]%[ma:d44c1a96]
	Một cấp số nhân có $6$ số hạng với công bội bằng $2$ và tổng số các số hạng  bằng $189$. Tìm số hạng cuối $u_6$ của cấp số nhân đã cho.
	 \shortans{96}
	\loigiai{ 
		Ta  có
		 \[
		\heva{ 
		&q = 2\\ 
		 &189 =S_6 = u_1 \cdot \dfrac{1 - q^6}{1 - q} = u_1 \cdot \dfrac{1 - 2^6}{1 - 2}
		} 
		 \Leftrightarrow
		\heva{ 
		 &q = 2\\
		&u_1 = 3 
		 }
		\Rightarrow u_6 = u_1 q^5 = 3 \cdot 2^5 =  96.
		 \]
	} 
\end{ex}

% nguồn: Ok_TraLoiNgan_Lop11_Lan2
\begin{ex}%[1D2H3-4]%[ma:f5ef7ea3]
	Cho cấp số nhân $(u_n)$ có công bội $q$ và  thỏa mãn
	 \[
	\heva{ 
	&u_1 + u_2 + u_3 + u_4 +  u_5 = 49\left(\dfrac{1}{u_1} + \dfrac{1}{u_2} + \dfrac{1}{u_3} + \dfrac{1}{u_4} + \dfrac{1}{u_5}\right)\\
	& u_1 + u_3 =  35.
	 }
	 \]
	Tính  $P = u_1 + 4q^2$.
	 \shortans{29}
	 \loigiai{
		Giả sử $(u_n)$ có công bội $q=0$.  Khi đó $u_2=u_3=u_4=u_5=0$ (vô lý).\\
		Giả sử $(u_n)$ có công bội $q=1$. Khi đó  $u_1=u_2=u_3=u_4=u_5$.\\
		 Khi đó $\heva{
		&5u_1  = 49\cdot \dfrac{5}{u_1} \\
		& 2u_1= 35 
		}$  (vô lý).\\
		Xét $q\ne 0$, $q\ne 1$.\\ 
		Nhận xét: Nếu $u_1$, $u_2$, $u_3$, $u_4$, $u_5$ là một cấp số nhân với công bội $q$  thì $\dfrac{1}{u_1}$, $\dfrac{1}{u_2}$, $\dfrac{1}{u_3}$, $\dfrac{1}{u_4}$, $\dfrac{1}{u_5}$ cũng tạo thành cấp số nhân với công bội $\dfrac{1}{q}$.\\
		Do  đó từ giả thiết ta có $\heva{& u_1 \cdot \dfrac{q^5 - 1}{q - 1} = 49 \cdot \left( \dfrac{1}{u_1} \cdot \dfrac{1 - \dfrac{1}{q^5}}{1 - \dfrac{1}{q}} \right) \quad(1)\\ &u_1 + u_1 q^2 = 35. \quad (2)}$\\
		Phương trình (1) $\Leftrightarrow u_1 \cdot \dfrac{q^5 -  1}{q - 1} = \dfrac{49}{u_1} \cdot \left( \dfrac{q^5 - 1}{q^4(q - 1)} \right) \Leftrightarrow u_1^2 q^4 = 49 \Leftrightarrow u_1 q^2 = \pm 7$.\\
		Với $u_1 q^2 = -7$. Thay vào (2), ta được $u_1 - 7 = 35 \Leftrightarrow u_1 =  42$. Suy ra $q^2 = \dfrac{-7}{42}$ (vô lý).\\
		Với $u_1 q^2 = 7$. Thay vào (2), ta được $u_1 +  7 = 35 \Leftrightarrow u_1 = 28$. \\
		Do  đó
		\[ 
		\heva{&u_1 = 28\\&q = \dfrac{1}{2}}  \quad \text{hoặc} \quad \heva{u_1 = 28\\q = -\dfrac{1}{2}.}
		\] 
		Khi đó $u_1 + 4q^2  = 29$.
	} 
\end{ex}

% nguồn: Ok_TraLoiNgan_Lop11_Lan2
\begin{ex}%[1D2H3-4]%[ma:fcd6c9b3]
	Một cấp số nhân có công bội bằng $3$ và số hạng đầu bằng $5$. Biết số hạng chính giữa là $32805$. Hỏi cấp  số nhân đã cho có bao nhiêu số hạng?
	 \shortans{17}
	 \loigiai{
		 \[
		32\,805 = u_n = u_1 q^{n-1} = 5 \cdot 3^{n-1}  \Leftrightarrow 3^{n-1} = 6\,561 = 3^8 \Leftrightarrow n = 9.
		\] 
		Vậy $u_9$ là số hạng chính giữa  của cấp số nhân, nên cấp số nhân đã cho có $17$ số hạng.
	} 
\end{ex}

% nguồn: Ok_TraLoiNgan_Lop11_Lan1
\begin{ex}%[1D2V3-4]%[ma:25d4ee76]
	Tìm công bội của cấp số nhân thỏa $\heva{&u_1+u_2+u_3=135\\ &u_4+u_5+u_6=40}$ là $\dfrac{a}{b}$ với $a$, $b$ là hai số tự nhiên nguyên  tố cùng nhau. Giá trị $a+b$
	\shortans{5} 
	 \loigiai{
		$\heva{&u_1+u_2+u_3=135\\ 
		 &u_4+u_5+u_6=40} \Leftrightarrow\heva{&u_1+u_1\cdot q+u_1\cdot q^2=135\\
		&u_1\cdot q^3+u_1\cdot q^4+u_1\cdot  q^5=40} \Leftrightarrow\heva{&u_1\left(1+q+q^2\right)=135 \quad (1)\\
		 &u_1\cdot q^3\left(1+q+q^2\right)=40 \quad (2)}$.\\
		Lấy (2) chia cho (1) ta được: $q^3=\dfrac{8}{27} \Rightarrow q=\dfrac{2}{3}  \Rightarrow u_1=\dfrac{1215}{19}$.
	 }
\end{ex}

% nguồn: Ok_TraLoiNgan_Lop11_Lan1
\begin{ex}%[1D2V3-4]%[ma:eaf7d891]
	Cho hình vuông $A_1$ có cạnh  bằng $ 1$. Gọi $A_2$ là hình vuông có các đỉnh là trung điểm các cạnh của hình vuông $A_1$; $A_3$ là hình vuông có các đỉnh là trung điểm các cạnh của hình vuông $A_2$,... Cứ tiếp tục quá trình như trên ta được dãy các hình vuông $A_1;\,A_2;\,...;\,A_n;...$ Biết rằng diện tích của hình vuông $A_{2024}$ là $S=\dfrac{1}{2^{\overline{abcd}}}$. Tính $a+b+c+d$.
	\shortans{7} 
	\loigiai{ 
		 Hình vuông $A_1$ có diện tích $S_1=1$.\\
		Hình vuông $A_2$ là hình vuông có các đỉnh là trung điểm  các cạnh của hình vuông $A_1$ do đó hình vuông $A_2$ có diện tích $S_2=\dfrac{1}{2}{S_1}=\dfrac{1}{2}$.\\
		Tương tự, hình vuông $A_3$ có diện tích  $S_3=\dfrac{1}{2}{S_2}=\dfrac{1}{2}\cdot\dfrac{1}{2}=\dfrac{1}{2^2}$.\\
		 Cứ tiếp tục như thế ta tính được diện tích hình vuông $A_{2024}$ là $S_{2024}=\dfrac{1}{2^{2023}}$.
	} 
\end{ex}

% nguồn: Ok_TraLoiNgan_Lop11_Lan1
\begin{ex}%[1D2V3-4]%[ma:509568a7]
	Tổng ba số hạng liên tiếp của một cấp số cộng là $21$. Nếu lấy số thứ hai trừ đi $1$ và số thứ ba cộng thêm $1$ thì ba số đó lập thành  một cấp số nhân. Tìm tổng ba số hạng đầu, biết số hạng đầu có giá trị nhỏ hơn $4$.
	\shortans{21} 
	\loigiai{ 
		 Gọi $\left(u_n\right)$ là cấp số cộng có công sai $d$.\\
		Ta có $u_1+u_2+u_3=21\Rightarrow u_1+\left(  u_1+d \right)+\left( u_1+2d \right)=21\Rightarrow u_1+d=7$\hfill{(*)}\\
		Theo giả thiết ${u_1;\, u_2-1;\, u_3+1}$ lập thành cấp số nhân hay $u_1;\,  u_1+d-1;\, u_1+2 d+1$ lập thành một cấp số nhân.\\
		Suy ra ${\left(u_1+d-1\right)^2=u_1\left(u_1+2 d+1\right)  \Rightarrow\left(u_1+d-1\right)^2=u_1\left[2\left(u_1+d\right)-u_1+1\right]}$.\\
		Kết hợp với  $(*)$, ta được
		$$ 
		\begin{aligned}
			&(7-1)^2=u_1\left(2.7-u_1+1\right)\\ 
			&  \Leftrightarrow 36=u_1\left(15-u_1\right)\\
			& \Leftrightarrow  u_1^2-15u_1+36=0\\
			&\Leftrightarrow \hoac{ 
			 &u_1=12>4\quad \text{(loại)} \\
			&u_1=3 
			} 
		\end{aligned}
		$$ 
		Với $u_1=3$ thì  $d=4\Rightarrow u_2=7;\, u_3=11$.\\
		Vậy  $u_1+u_2+u_3=3+7+11=21$.
	 }
\end{ex}

% nguồn: Ok_TraLoiNgan_Lop11_Lan2
\begin{ex}%[1D2V3-4]%[ma:512b79e7]
	Cho cấp số cộng $(u_n)$ gồm bốn  số nguyên. Biết tổng của chúng bằng $20$, tổng các nghịch đảo của chúng bằng $\dfrac{25}{24}$. Tìm tổng $4$ số đó.
	 \shortans{20}
	\loigiai{ 
		Gọi $d = 2x$ là công sai của cấp số cộng đã cho, vì  cấp số cộng đã cho gồm bốn số nguyên nên suy ra $d \in \mathbb{Z}$.\\
		Ta có bốn số cần tìm là  $a - 3x$; $a - x$; $a + x$; $a + 3x$.
		\begin{eqnarray*}
			&&(a - 3x) + (a - x) +  (a + x) + (a + 3x) = 20 \\
			&\Leftrightarrow& 4a  = 20\\
			 &\Leftrightarrow& a = 5.
		\end{eqnarray*}
		Ta  lại có
		\begin{eqnarray*}
			\dfrac{1}{a  - 3x} + \dfrac{1}{a - x} + \dfrac{1}{a + x} + \frac{1}{a + 3x}
			 &=& \dfrac{25}{24} \\
			\Leftrightarrow \dfrac{1}{5 - 3x} + \dfrac{1}{5  - x} + \dfrac{1}{5 + x} + \dfrac{1}{5 + 3x}
			&=& \dfrac{25}{24} \quad  (1)
		\end{eqnarray*}
		Điều  kiện
		 \[
		 \heva{
		& x \neq  -\frac{5}{3} \\
		 & x \neq -5 \\
		& x \neq \frac{5}{3}  \\
		& x \neq 5. 
		} 
		 \]
		Giải $(1)$ 
		\begin{eqnarray*}
			&&\frac{1}{5 - 3x} + \frac{1}{5 + 3x}  + \frac{1}{5 - x} + \frac{1}{5 + x} = \frac{25}{24} \\
			&\Leftrightarrow& \frac{10}{25 - 9x^2} +  \frac{10}{25 - x^2} = \frac{25}{24} \\
			&\Leftrightarrow& \frac{2}{25 - 9x^2} + \frac{2}{25 -  x^2} = \frac{5}{24} \\
			&\Leftrightarrow& 48(25 -  x^2) + 48(25 - 9x^2) = 5(25 - x^2)(25 - 9x^2) \\
			&\Leftrightarrow& 45x^4 -  770x^2 + 725 = 0\\
			 &\Leftrightarrow& \hoac{
			& x^2 = 1 \Rightarrow x =  \pm 1 \\
			 & x^2 = \frac{145}{9} \Rightarrow x = \pm \frac{\sqrt{145}}{3}
			} 
		\end{eqnarray*}
		Đối chiếu  điều kiện và vì $d \in \mathbb{Z}$ nên ta nhận $x = \pm 1$.\\
		Với $a = 5, x = 1$: ta có $u_1 = 2, u_2 = 4, u_3 = 6, u_4 = 8$.\\ 
		Với $a = 5, x = -1$: ta có $u_1 = 8, u_2 = 6, u_3 = 4, u_4 =  2$.\\
		Vậy tổng của $4$  chữ số là $2+4+6+8=20$.
	} 
\end{ex}

% nguồn: Ok_TraLoiNgan_Lop11_Lan2
\begin{ex}%[1D2V3-4]%[ma:94bb9898]
	Tìm số ở giữa của ba số khác nhau biết chúng tạo thành cấp số cộng có tổng bằng $6$, nếu hoán đổi vị trí số hạng thứ nhất và số  hạng thứ hai đồng thời giữ nguyên số hạng thứ ba ta được cấp số nhân.
	\shortans{2} 
	\loigiai{ 
		Gọi ba số cần tìm là $u_1$, $u_2$, $u_3$ với $u_1$, $u_2$, $u_3$ đôi một  khác nhau.\\
		Vì $u_1$, $u_2$, $u_3$ tạo  thành cấp số cộng với công sai $d \neq 0$ nên $u_1$, $u_2=u_1+d$, $u_3=u_1+2 d$.\\
		Hơn nữa, $u_1+u_2+u_3=6 \Leftrightarrow u_1+\left(u_1+d\right)+\left(u_1+2 d\right)=6 \Leftrightarrow u_1+d=2$.\\ 
		Vì $u_2$, $u_1$, $u_3$ tạo thành cấp số  nhân nên $u_1+d$, $u_1$, $u_1+2 d$ tạo thành cấp số nhân.\\
		Lúc  đó
		\begin{eqnarray*}
			\left(u_1+d\right)\left(u_1+2  d\right)&=&u_1^2
			\Leftrightarrow\left(u_1+d\right)\left(u_1+d+d\right)=u_1^2 \Leftrightarrow 2\left(2+2-u_1\right)\\ 
			&=&u_1^2 \Leftrightarrow u_1^2+2 u_1-8=0 \Leftrightarrow\hoac{&u_1=2 \\& u_1=-4.} 
		\end{eqnarray*}
		Với $u_1=2$, suy ra $d=0$ không thỏa mãn.\\ 
		Với $u_1=-4$,  suy ra $d=6$.\\
		 Vậy ba số cần tìm là $-4$, $2$, $8$.
	} 
\end{ex}

% nguồn: Ok_TraLoiNgan_Lop11_Lan2
\begin{ex}%[1D2V3-4]%[ma:fd819772]
	Ba số hạng đầu của một cấp số nhân là $x-6$; $x$ và $y$. Tìm $y$,  biết rằng công bội của cấp số nhân là $6$ (kết quả làm tròn đến hàng đơn vị).
	 \shortans{43,2}
	 \loigiai{
		Cấp số  nhân có ba số hạng đầu là $x-6$; $x$ và $y$ có công bội $q=6$, nên ta có
		\[\heva{&u_1=x-6 ,\quad  q=6\\&x=u_2=u_1q=6(x-6)\\&y=u_3=u_1q^2=36(x-6)} \Rightarrow \heva{&x=\dfrac{36}{5}\\&
		y=36\cdot \left(\dfrac{36}{5}-6\right)= 43{,}2.}\] 
	} 
\end{ex}

% nguồn: Ok_TraLoiNgan_Lop11_Lan2
\begin{ex}%[1D2V3-4]%[ma:36a0a1d4]
	Một tam giác có ba cạnh lập thành một cấp số nhân. Biết cạnh lớn nhất hơn cạnh nhỏ nhất $10$ cm và độ đài $3$ cạnh đều là số tự nhiên. Tính  tổng độ dài ba cạnh của tam giác đã cho.
	\shortans{38} 
	\loigiai{ 
		Gọi độ dài ba cạnh của  tam giác là $a$, $b$, $c$ $\mathrm{cm}$.\\
		Không mất tính tổng quát giả sử $0<c<b<a$.\\ 
		Theo  bài ra, ta có
		 \[
		 \heva{&
		 a=c+10 \\&
		b^2=a c  \\&
		 a, b, c \in \mathbb{N}.}
		 \]
		 Do đó
		\[ 
		b^2=c(c+10) \Leftrightarrow  c^2+10 c+25-b^2=25 \Leftrightarrow(c+5-b)(c+5+b)=25.
		\] 
		Suy ra $(c+5+b) \in \{ \pm 1; \pm 5; \pm 25\}$ mà $c+5+b>c+5-b$, $c+5+b>0$ nên ta  chỉ có một trường hợp là
		\[\heva{&c+5+b= 25\\&c+5-b=1} \Leftrightarrow \heva{&c+b= 20\\&c-b=- 4} \Leftrightarrow \heva{& c=8  \\& b=12.}
		 \]
		Từ  đó có $a=c+10=18$.\\
		 Vậy $a+b+c=38$.
	} 
\end{ex}

% nguồn: Ok_TraLoiNgan_Lop11_Lan2
\begin{ex}%[1D2V3-4]%[ma:e52b113d]
	Bốn góc của một tứ giác tạo thành cấp số nhân và góc lớn nhất gấp  $27$ lần góc nhỏ nhất. Tổng của góc lớn nhất và góc bé nhất bằng (đơn vị độ) bằng bao nhiêu?
	\shortans{252} 
	\loigiai{ 
		Giá sử bốn góc $A$,  $B$, $C$, $D$ theo thứ tự đó lập thành cấp số nhân thóa yêu cầu với công bội $q$. \\
		Ta  có
		\[ 
		\heva{&A+B+C+D=  360^\circ\\&D=27A} \Leftrightarrow \heva{&A(1+q+q^2+q^3)=360^\circ\\&Aq^3 = 27A}\Leftrightarrow \heva{&
		q=3 \\&A=9^\circ\\& 
		D=A q^3=243^\circ}\Rightarrow  A+D=252^\circ.
		 \]
	} 
\end{ex}

% nguồn: Ok_TraLoiNgan_Lop11_Lan2
\begin{ex}%[1D2V3-4]%[ma:195b42a3]
	Số đo ba kích thước của hình hộp chữ nhật lập thành  một cấp số nhân. Biết thể tích của khối hộp là $3\,375 \mathrm{~cm}^3$ và diện tích toàn phần là $1\,950 \mathrm{~cm}^2$. Tính tổng số đo ba kích thước của hình hộp chữ nhật đó.
	 \shortans{65}
	 \loigiai{
		Gọi ba kích thước của hình hộp lập thành một cấp số nhân lần lượt là $x$, $xq$,  $xq^2$ $(\mathrm{cm}), x, q>0$.\\
		Vì thể tích của khối hộp là $3\,375 \mathrm{~cm}^3$ nên $x  \cdot q q^2=3\,375 \Leftrightarrow x^3 \cdot q^3=3\,375 \Leftrightarrow x q=15$.\\
		Vì diện  tích toàn phần là $1\,950 \mathrm{~cm}^2$ nên
		\begin{eqnarray*}
			& &2 x q \cdot x q^2+2 x  \cdot x q+2 x \cdot x q^2=1\,950 \\
			&\Leftrightarrow& x^2 q^2 \cdot q+x \cdot x q+x^2 q^2=975 \\ 
			 & \Leftrightarrow& 15^2 \cdot q+x \cdot 15+15^2=975 \\
			 & \Leftrightarrow &225 q+15 x=750 \Leftrightarrow x=50-15 q
		\end{eqnarray*}
		Thay $x=50-15 q$ vào $x  q=15$, ta được $(50-15 q) \cdot q=15 \Leftrightarrow\hoac{&q=3 \\&q=\dfrac{1}{3}.}$\\
		Với $q=3 \Rightarrow x=5 \Rightarrow$ độ dài ba  kích thước là $5$, $15$, $45$.\\
		Với $q=\dfrac{1}{3} \Rightarrow x=45 \Rightarrow$  độ ba kích thước là $45$, $15$, $5$.\\
		Vậy tổng độ dài ba kích thước là  $45+15+5=65$.
	} 
\end{ex}

% nguồn: Ok_TraLoiNgan_Lop11_Lan1
\begin{ex}%[1D2C3-4]%[ma:d8c75e9a]
	Cho cấp số nhân $(u_n)$ có $u_1=3$ và $15 u_1-4 u_2+u_3$ đạt giá trị nhỏ nhất. Tìm số hạng thứ $10$ của  cấp số nhân đã cho.
	 \shortans{1536}
	 \loigiai{
		Gọi ${q}$ là công bội của cấp số nhân ${(u_n)}$. Ta có: ${u_2=u_1 q=3 q; u_3=u_1  q^2=3 q^2}$.\\
		Suy  ra $15 u_1-4 u_2+u_3=45-12 q+3 q^2=3(q-2)^2+33 \geq 33, \forall q \in \mathbb{R}$.\\
		 Ta có: $15 u_1-4 u_2+u_3$ đạt giá trị nhỏ nhất (bằng 33) khi và chỉ khi ${q=2}$.\\
		Khi đó: $u_{10}=u_1 q^{9}=3 \cdot 2^{9}=1536$. 
	} 
\end{ex}


%%%%%%%%%%%%%%%% TỰ LUẬN %%%%%%%%%%%%%%%%

% nguồn: Ok_TuanLuan_Toan11_Lan2
\begin{ex}%[1D2H3-4]%[ma:7a80bdcc]
	Tìm số hạng  đầu và công bội của cấp số nhân, biết
	\begin{multicols}{2}
		\begin{enumerate}
			\item  $\heva{&u_1+u_5=51\\&u_2+u_6=102.}$
			\item $\heva{&u_1+u_2+u_3=135\\&u_4+u_5+u_6=40.}$ 
		\end{enumerate}
	\end{multicols}
	 \loigiai{
		\begin{enumerate}
			\item  $\heva{&u_1+u_5=51\\&u_2+u_6=102}
			\Leftrightarrow \heva{&u_1+u_1q^4=51\\&u_1q+u_1q^5=102} 
			 \Leftrightarrow \heva{&u_1 \left(1+q^4\right)=51~\quad(\ast)\\&u_1q\left(1+q^4\right)=102~\quad(\ast\ast)}$\\
			 Ta có
			 \[
			 \dfrac{(\ast)}{(\ast\ast)} \Leftrightarrow \dfrac{u_1q(1+q^4)}{u_1(1+q^4)} = \dfrac{102}{51}
			\Leftrightarrow q=2 \Rightarrow u_1 = \dfrac{51}{1+q^4} = \dfrac{51}{17} =  3.
			\] 
			Kết  luận công bội $q=2$ và số hạng đầu tiên $u_1=3$.
			\item  $\heva{&u_1+u_2+u_3=135\\&u_4+u_5+u_6=40}
			\Leftrightarrow \heva{&u_1+u_1q+u_1q^2=135\\&u_1q^3+u_1q^4+u_1q^5=40} 
			\Leftrightarrow \heva{&u_1(1+q+q^2)=135~\quad(\ast)\\&u_1q^3(1+q+q^2)=40~\quad(\ast\ast)}$\\ 
			Ta có 
			\[ 
			\dfrac{(\ast\ast)}{(\ast)}  \Leftrightarrow \dfrac{u_1q^3(1+q+q^2)}{u_1(1+q+q^2)} = \dfrac{40}{135}
			\Leftrightarrow q^3 = \dfrac{8}{27} \Leftrightarrow  q = \sqrt[3]{\dfrac{8}{27}} = \dfrac{2}{3}.
			 \]
			Thay $q = \dfrac{2}{3}$ vào phương  trình $(\ast)$
			\begin{eqnarray*}
				&  &u_1\left(1+\dfrac{2}{3}+\left(\dfrac{2}{3}\right)^2\right) = 135\\
				&\Leftrightarrow &u_1\left(1+\dfrac{2}{3}+\dfrac{4}{9}\right)  = 135\\
				&\Leftrightarrow  &u_1\left(\dfrac{9+6+4}{9}\right) = 135\\
				&\Leftrightarrow  &u_1\left(\dfrac{19}{9}\right) = 135\\
				&\Leftrightarrow &u_1  = \dfrac{135 \cdot 9}{19} = \dfrac{1215}{19}.
			\end{eqnarray*}
			Kết luận công bội $q = \dfrac{2}{3}$ và số hạng  đầu tiên $u_1=\dfrac{1215}{19}$.
		\end{enumerate}
	} 
\end{ex}

% nguồn: Ok_TuanLuan_Toan11_Lan2
\begin{ex}%[1D2H3-4]%[ma:ece70a05]
	Tìm số  hạng đầu và công bội của cấp số nhân
	\begin{multicols}{3}
		\begin{enumerate}[a)]
			\item $\heva{ 
			 &u_2+u_3+u_4=\dfrac{35}{2}\\
			&u_1u_5=25\\ 
			&u_i>0~(i=1,\cdots  5);}$
			 \item $\heva{
			&u_1-u_3+u_5=65\\ 
			&u_1+u_7=325;}$ 
			\item 
			 $\heva{
			&u_2+u_4+u_6=-42\\ 
			&u_3+u_5=20;}$ 
			 \item $\heva{
			&u_1+u_6=165\\ 
			 &u_3+u_4=60;}$
			\item  $\heva{
			&u_1+u_2+u_3+u_4=15\\ 
			&u_1^2+u_2^2+u_3^2+u_4^2=85;}$ 
			 \item $\heva{
			 &u_1+u_2+u_3=13\\
			 &u_4+u_5+u_6=351.}$
		\end{enumerate}
	\end{multicols}
	\loigiai{ 
		\begin{enumerate}
			 \item
			 $\heva{
			 &u_2+u_3+u_4=\dfrac{35}{2}\\
			&u_1u_5=25\\ 
			 &u_i>0~(i=1, \dots, 5)
			}  \Leftrightarrow \heva{
			&u_1q+u_1q^2+u_1q^3 = \dfrac{35}{2}~\quad(1)\\ 
			 &u_1u_1q^4 = 25~\quad(2)
			}$\\ 
			Ta có 
			$(2) \Leftrightarrow (u_1q^2)^2=25 \Leftrightarrow u_1q^2=5 \Leftrightarrow u_1=\dfrac{5}{q^2}$ thay vào $(1)$  được
			 \[
			\dfrac{5}{q^2}(q+q^2+q^3) = \dfrac{35}{2}  \Leftrightarrow 2(1+q+q^2) = 7q \Leftrightarrow 2q^2-5q+2=0 \Leftrightarrow \hoac{&q=2\\&q=\dfrac{1}{2}.}
			\] 
			Với $q=2 \Rightarrow u_1=\dfrac{5}{4}$; với $q=\dfrac{1}{2} \Rightarrow u_1=20$. 
			\item 
			 $\heva{&u_1-u_3+u_5=65\\&u_1+u_7=325.}
			\Leftrightarrow \heva{&u_1-u_1q^2+u_1q^4=65\\&u_1+u_1q^6=325.} 
			 \Leftrightarrow \heva{
			&u_1(1-q^2+q^4)=65~\quad(1)\\ 
			 &u_1(1+q^6)=325~\quad(2).}$
			\begin{eqnarray*}
				&\dfrac{(2)}{(1)} &\Leftrightarrow \dfrac{1+q^6}{1-q^2+q^4}=\dfrac{325}{65}\\ 
				&  &\Leftrightarrow \dfrac{(1+q^2)(1-q^2+q^4)}{1-q^2+q^4}=5\\
				& &\Leftrightarrow  1+q^2=5\\
				& &\Leftrightarrow  q^2=4 \Leftrightarrow q=\pm 2.
			\end{eqnarray*}
			Với  $q=2 \Rightarrow u_1=\dfrac{65}{1-2^2+2^4}=5$;
			với  $q=-2 \Rightarrow u_1=\dfrac{65}{1-(-2)^2+(-2)^4}=5$.
			\item  $\heva{&u_2+u_4+u_6=-42\\&u_3+u_5=20}
			 \Leftrightarrow \heva{&u_1q+u_1q^3+u_1q^5=-42\\&u_1q^2+u_1q^4=20}
			 \Leftrightarrow \heva{&u_1q (1+q^2+q^4)=-42~\quad(1)\\&u_1q^2(1+q^2)=20~\quad(2)}$\\
			 Lấy
			\begin{eqnarray*}
				&\dfrac{(1)}{(2)}&\Leftrightarrow  \dfrac{1+q^2+q^4}{q(1+q^2)} = -\dfrac{21}{10}\\
				 & &\Leftrightarrow 10+10q^2+10q^4=-21q-21q^3\\
				& &\Leftrightarrow 10q^4+21q^3+10q^2+21q+10=0\\ 
				&  &\Leftrightarrow 10q^2+21q+10+\dfrac{21}{q}+\dfrac{10}{q^2}=10\\
				&  &\Leftrightarrow 10\left(q^2+\dfrac{1}{q^2}\right)+21\left(1+\dfrac{1}{q}\right)+10=0~\quad(\ast).
			\end{eqnarray*}
			Đặt $t=q+\dfrac{1}{q} \Rightarrow t^2=\left(q+\dfrac{1}{q}\right)^2 \Leftrightarrow q^2+\dfrac{1}{q^2}=t^2-2$. Điều kiện $|t| \geq 2$.\\ 
			$(\ast) \Leftrightarrow 10(t^2-2)+21t+10=0  \Leftrightarrow 10t^2+21t-10=0 \Leftrightarrow \hoac{
			&t = -\dfrac{5}{2}~(\text{nhận})\\ 
			 &t = \dfrac{2}{5}~(\text{loại})
			}$\\ 
			Với $t=-\dfrac{5}{2} \implies q+\dfrac{1}{q}=-\dfrac{5}{2} \Leftrightarrow  2q^2+5q+2=0 \Leftrightarrow \hoac{
			 &q=-\dfrac{1}{2}\\
			&q=-2}$\\ 
			Nếu $q=-\dfrac{1}{2} \implies u_1 = \dfrac{20}{q^2+q^4}  = \dfrac{20}{\left(-\dfrac{1}{2}\right)^2+\left(-\dfrac{1}{2}\right)^4}=64$.\\
			Nếu  $q=-2 \implies u_1 = \dfrac{20}{q^2+q^4} = \dfrac{20}{(-2)^2+(-2)^4}=1$.
			 \item $\heva{&u_1+u_6=165\\&u_3+u_4=60}
			\Leftrightarrow \heva{&u_1(1+q^5)=165~\quad(1)\\&u_1q^2(1+q)=60~\quad(2)}$ 
			\begin{eqnarray*}
				&\dfrac{(1)}{(2)} &\Leftrightarrow \dfrac{1+q^5}{q^2(1+q)} = \dfrac{165}{60}  = \dfrac{11}{4}\\
				& &\Rightarrow  4(1-q+q^2-q^3+q^4)=11q^2\\
				 & &\Leftrightarrow 4q^4-4q^3-7q^2-4q+4=0\\
				&  &\Leftrightarrow 4\left(q^2+\dfrac{1}{q^2}\right)-4\left(q+\dfrac{1}{q}\right)-7=0~\quad(\ast)
			\end{eqnarray*}
			Đặt  $t=q+\dfrac{1}{q} \Rightarrow t^2=\left(q+\dfrac{1}{q}\right)^2 \Rightarrow q^2+\dfrac{1}{q^2}=t^2-2$. Điều kiện: $|t|\ge 2$.\\
			$(\ast) \Leftrightarrow 4(t^2-2)-4t-7=0 \Leftrightarrow  4t^2-4t-15=0 \Leftrightarrow \hoac{&t=\dfrac{5}{2}\\&t=-\dfrac{3}{2}~(\text{loại})}$\\
			Với $t=\dfrac{5}{2}  \Rightarrow q+\dfrac{1}{q}=\dfrac{5}{2} \Leftrightarrow 2q^2-5q+2=0
			\Leftrightarrow  \hoac{&q=2\\&q=\dfrac{1}{2}.}$\\
			Với $q=2 \Rightarrow  u_1=\dfrac{165}{1+q^5}=\dfrac{165}{1+2^5}=5$.\\
			Với $q=\dfrac{1}{2} \Rightarrow u_1=\dfrac{165}{1+q^5}=\dfrac{165}{1+\left(\dfrac{1}{2}\right)^5}=160$. 
			\item 
			Từ  $\heva{&u_1+u_2+u_3+u_4=15\\&u_1^2+u_2^2+u_3^2+u_4^2=85}
			 \Leftrightarrow \heva{&u_1(1+q+q^2+q^3)=15~\quad(1)\\&u_1^2(1+q^2+q^4+q^6)=85~\quad(2).}$
			 \\
			 Ta có
			\begin{eqnarray*}
				&\dfrac{(1)}{(2)} &\Leftrightarrow \dfrac{u_1(1+q+q^2+q^3)}{u_1^2(1+q^2+q^4+q^6)} =  \dfrac{15}{85}\\
				& &\Rightarrow  \dfrac{1+q+q^2+q^3}{u_1(1+q^2+q^4+q^6)}=\dfrac{3}{17}\\
				 & &\Leftrightarrow \dfrac{(1+q)(1+q^2)}{u_1(1+q^2)(1+q^4)} = \dfrac{3}{17}\\
				& &\Rightarrow \dfrac{1+q}{u_1(1+q^4)} = \dfrac{3}{17}.\\ 
				& &\Rightarrow  17(1+q)=3u_1(1+q^4)\\
				&  &\Leftrightarrow u_1=\dfrac{17(1+q)}{3(1+q^4)}.
			\end{eqnarray*}
			Thay  $u_1$ vào $(1)$ được
			\begin{eqnarray*}
				 & &\dfrac{17(1+q)}{3(1+q^4)}(1+q+q^2+q^3)=15\\
				&\Leftrightarrow &\dfrac{17(1+q)^2(1+q^2)}{3(1+q^4)}=15\\ 
				&\Leftrightarrow  &17(1+2q+q^2)(1+q^2) = 45(1+q^4)\\
				&\Leftrightarrow  &17(1+2q+2q^2+2q^3+q^4)=45(1+q^4)\\
				 &\Leftrightarrow &17+34q+34q^2+34q^3+17q^4=45+45q^4\\
				&\Leftrightarrow  &28q^4-34q^3-34q^2-34q+28=0.\\
			\end{eqnarray*}
			Vì $q=0$ không phải là nghiệm, chia hai  vế cho $q^2$:
			\[ 
			28q^2-34q-34-\dfrac{34}{q}+\dfrac{28}{q^2}=0 \Leftrightarrow 14\left(q^2+\dfrac{1}{q^2}\right)-17\left(q+\dfrac{1}{q}\right)-17=0~\quad(\ast) 
			 \]
			Đặt  $t=q+\dfrac{1}{q} \Rightarrow t^2=\left(q+\dfrac{1}{q}\right)^2 \Rightarrow q^2+\dfrac{1}{q^2}=t^2-2$. Điều kiện: $|t|\ge 2$.\\
			$(\ast)  \Leftrightarrow 14(t^2-2)-17t-17=0 \Leftrightarrow 14t^2-17t-45=0 \Leftrightarrow \hoac{
			&t=\dfrac{5}{2}\\ 
			 &t=-\dfrac{9}{7}~(\text{loại}).
			 }$
			\begin{itemize}
				\item Với $t=\dfrac{5}{2}  \Rightarrow q+\dfrac{1}{q}=\dfrac{5}{2} \Leftrightarrow 2q^2-5q+2=0 \Leftrightarrow \hoac{&q=2\\&q=\dfrac{1}{2}.}$
				\item Với $q=2 \Rightarrow u_1=1$; với $q=\dfrac{1}{2} \Rightarrow  u_1=\dfrac{15}{1+q+q^2+q^3}=8.$
			\end{itemize}
		\end{enumerate}
	} 
\end{ex}

% nguồn: Ok_TuanLuan_Toan11_Lan2
\begin{ex}%[1D2H3-4]%[ma:85cbe789]
	Cho cấp số nhân $(u_n)$  có các số hạng thỏa mãn:
	 $\heva{&u_1+u_5=51\\&u_2+u_6=102.}$
	\begin{enumerate}[a)]
		\item Tìm số hạng  đầu và công bội của cấp số nhân.
		\item Số $12\,288$ là số hạng thứ mấy? 
	\end{enumerate}
	\loigiai{ 
		\begin{enumerate}
			\item Ta có: 
			 $\heva{&u_1+u_5=51\\&u_2+u_6=102}
			\Leftrightarrow  \heva{&u_1+u_1q^4=51\\&u_1q+u_1q^5=102}
			 \Leftrightarrow \heva{&u_1 \left(1+q^4\right)=51~\quad(\ast)\\&u_1q\left(1+q^4\right)=102~\quad(\ast\ast)}$\\
			Ta  có
			\[ 
			\dfrac{(\ast\ast)}{(\ast)} \Leftrightarrow \dfrac{u_1q(1+q^4)}{u_1(1+q^4)} =  \dfrac{102}{51}
			\Leftrightarrow q=2 \Rightarrow u_1 = \dfrac{51}{1+q^4} = \dfrac{51}{1+2^4} = \dfrac{51}{17} =  3.
			 \]
			Kết luận công bội  $q=2$ và số hạng đầu tiên $u_1=3$.
			\item Ta  có $u_k=12288 \Leftrightarrow u_1q^{k-1}=12\,288$.\\
			Thay $u_1=3$ và $q=2$ vào,  ta được:
			$3 \cdot 2^{k-1} = 12\,288 \Leftrightarrow 2^{k-1} = \dfrac{12\,288}{3}  = 4\,096$.\\
			Vì $4\,096 = 2^{12}$, nên ta có $2^{k-1} = 2^{12} \Leftrightarrow k-1=12 \Leftrightarrow k=13$.\\ 
			Kết luận số $12\,288$ là số hạng thứ  $13$.
		\end{enumerate}
	} 
\end{ex}
